% ======================================================================
%  Constrained-worldline brachistochrones in non-stationary spacetimes
%  RevTeX 4.2 (PRD) skeleton -- fill prose marked [[...]]
% ======================================================================
\documentclass[
    aps, prd, twocolumn, superscriptaddress,
    showpacs, nofootinbib, floatfix
]{revtex4-2}

\usepackage{amsmath,amssymb,amsfonts}
\usepackage{graphicx}
\usepackage{bm}
\usepackage{xcolor}
\usepackage{hyperref}
\hypersetup{colorlinks=true, linkcolor=blue, citecolor=blue, urlcolor=blue}

% all figures collected here as PNG
\graphicspath{{Immagini/}}

% float placement (skeleton has many figures / sparse prose)
\setcounter{topnumber}{4}
\setcounter{bottomnumber}{4}
\setcounter{totalnumber}{10}
\renewcommand{\topfraction}{0.95}
\renewcommand{\bottomfraction}{0.95}
\renewcommand{\textfraction}{0.05}
\renewcommand{\floatpagefraction}{0.6}

% ToC: widen the number field so roman numerals don't collide with titles
\makeatletter
\renewcommand*\l@section{\@dottedtocline{1}{0em}{2.6em}}
\renewcommand*\l@subsection{\@dottedtocline{2}{2.6em}{2.4em}}
\makeatother

% prose placeholder helper
\newcommand{\todo}[1]{{\color{red}\small[[#1]]}}

% shortcuts
\newcommand{\Ehat}{\hat E}
\newcommand{\rp}{r_{+}}
\newcommand{\re}{r_{e}}
\newcommand{\Jc}{J_{c}}
\newcommand{\Af}{A_{\mathrm{freeze}}}
\newcommand{\dd}{\mathrm{d}}

\begin{document}

\title{Constrained-worldline brachistochrones in non-stationary
spacetimes: Kodama rails, closed forms, and plunge inversion}

\author{Iman Rosignoli}
\affiliation{University of Pavia, Pavia, Italy\\
\normalfont\texttt{iman.rosignoli01@universitadipavia.it}}

\date{\today}

\begin{abstract}
We extend the relativistic brachistochrone (fastest-descent
constrained rail curves, Perlick 1991) from stationary to non-stationary
spacetimes. The rail constraint $-u\cdot W=\Ehat$ is an actively controlled
invariant, legitimate via Pontryagin's principle; the selector $W$ follows a
hierarchy Killing $\to$ conformal Killing $\to$ Kodama vector. We treat FLRW
(degenerate base), Vaidya (dynamical black hole, $W=\partial_v$ the Kodama
vector), and Thakurta--Kerr (conformal rotating compact object). Main results:
Kodama energy conservation along the rail; the plunge-radius law and the
absence of a \emph{physical} evaporative inversion (only rotation and the
conformal factor invert); closed-form equatorial separatrices in Weierstrass
functions with Doran continuation through the ergosphere; the trichotomy at
the ergosphere with an analytic cusp theorem; an existence theorem for the
conformal plunge inversion; and closed-form first-order adiabatic corrections
whose transcendence is a universal genus-two hyperelliptic dilogarithm on the
frozen spectral curve, unified across the three geometries by a
breathing-indicatrix theorem in which the weight-two letter migrates from a
third-kind dipole at infinity (Thakurta--Kerr) to a $\mathrm{d}\!\log$ at the
horizon (Vaidya, where it is the face of the accretion/evaporation asymmetry).
\end{abstract}

\pacs{04.20.-q, 04.70.Bw, 98.80.Jk}
\maketitle

{\hypersetup{linkcolor=black}\tableofcontents}


% ======================================================================
\section{Introduction}
\label{sec:intro}
% ----------------------------------------------------------------------
The relativistic brachistochrone---the fastest constrained ``rail'' worldline
between two events---was formulated by Perlick for stationary spacetimes
\cite{Perlick1991,PerlickRayOptics}, where a timelike Killing vector provides
the conserved rail energy. It connects naturally to Fermat's principle and to
optical (Randers--Zermelo--Finsler) geometry
\cite{Kovner1990,Randers1941,GibbonsWerner2008,GibbonsRanders2009,BaoRoblesShen}.
Building on our earlier relativistic-brachistochrone studies---the stationary
Kerr geometry \cite{RosignoliKerr}, the interior Schwarzschild metric
\cite{RosignoliIntSchw}, and Tolman--Oppenheimer--Volkoff stellar interiors
\cite{RosignoliTOV}---we extend the construction to \emph{non-stationary}
spacetimes: FLRW, Vaidya \cite{Vaidya1951}, and Thakurta--Kerr
\cite{Thakurta1981}.
These extensions raise three questions that organize the paper.
\emph{(i) Is the rail constraint legitimate without a Killing vector?} We show it
is: $-u\cdot W=\Ehat$ is an actively controlled invariant, and the fastest
constrained worldline solves a well-posed time-optimal control problem
(Sec.~\ref{sec:principle}) whose control cost measures precisely the missing
symmetry. \emph{(ii) What replaces the conserved energy?} In a spherically
symmetric dynamical spacetime the selector $W$ is the Kodama vector, whose flux
charge is the Misner--Sharp mass; the rail conserves the Kodama energy
identically even though it is \emph{not} Noether-conserved for free fall
(Sec.~\ref{sec:vaidya}). \emph{(iii) Which phenomena are genuinely
non-stationary?} We isolate three: capture by cosmological expansion, the
teleological (future-anticipating) turning point in Vaidya, and the breathing
freezing surface that eats the radius of convergence of the quasi-constants in
Thakurta--Kerr. The three spacetimes form a ladder of increasing structure---%
FLRW (degenerate base), Vaidya (dynamical horizon), Thakurta--Kerr
(rotation $+$ conformal expansion)---treated in
Secs.~\ref{sec:flrw}--\ref{sec:thakurta}; equatorial closed forms and the
plunge-inversion analysis follow in Secs.~\ref{sec:closed}--\ref{sec:bvp}.


% ======================================================================
\section{The controlled-rail principle}
\label{sec:principle}
% ----------------------------------------------------------------------

\subsection{Rail constraint and the indicatrix}
\label{sec:indicatrix}
The rail is a constrained worldline with the invariant
\begin{equation}
    -\,u\cdot W \;=\; \Ehat ,
    \label{eq:rail}
\end{equation}
actively maintained by an ideal (magnetic-type) rail force $f_\mu=Du_\mu/\dd\tau$
with $f\cdot u=0$, so that the rail does no work on the particle.

\paragraph{The ellipse.} Fix an event and a time function adapted to $W$. The
linear constraint~\eqref{eq:rail} fixes the $W$-component of $u$; the mass shell
$u\cdot u=-1$ is a quadric; their intersection, projected onto the physical
velocity plane $(\dd r,\,r\,\dd\varphi)$ per unit selector-time, is an ellipse
\begin{equation}
    \dot x(\theta)=c+R(\theta),\qquad \theta\in[0,2\pi),
    \label{eq:ellipse}
\end{equation}
whose center $c$ is the Zermelo wind (zero for FLRW, ingoing radial for Vaidya,
azimuthal---frame dragging---for Thakurta--Kerr) and whose axes $R$ shrink with
the local speed as the freezing surface is approached (Fig.~\ref{fig:indicatrici}).
Every branch---arrival time $t$, proper time $\tau$, conformal time $\eta$---shares
this \emph{same} indicatrix; the branches differ only in the cost one-form, i.e.\
in which clock measures the elapsed travel time.

\paragraph{Pontryagin reduction.} Minimizing travel time subject
to~\eqref{eq:ellipse} is a time-optimal control problem with control $\theta$.
Since the admissible-velocity set is a smooth strictly convex oval, the maximizer
of $p\cdot\dot x(\theta)$ is unique and Pontryagin's maximum principle
\cite{Pontryagin,Liberzon} delivers the branch Hamiltonian $H(x,p)$ in closed
form (Secs.~\ref{sec:vaidya}--\ref{sec:thakurta}); the convex-oval regularity
makes the vakonomic and d'Alembert formulations coincide, so no Filippov
relaxation is needed. Two structural facts hold throughout: the transversality
condition
\begin{equation}
    H=0 \qquad\text{(free arrival time),}
    \label{eq:transversality}
\end{equation}
and $p_\varphi=J$ conserved (axial symmetry). In the stationary limit $H$ is
autonomous and $H\equiv0$ is a genuine first integral; non-stationarity enters
only through the explicit time dependence $\dd H/\dd(\text{time})=
\partial_{\text{time}}H$, proportional to $m'(v)$ (Vaidya) or $A'(\eta)$
(Thakurta--Kerr).

\paragraph{Cost of control.} Differentiating the invariant along the worldline,
\begin{equation}
    \frac{\dd(-u\cdot W)}{\dd\tau}=-\,u^a u^b\,\nabla_{(a}W_{b)},
    \label{eq:cost}
\end{equation}
which vanishes for \emph{all} $u$ iff $W$ satisfies the Killing equation. When it
does not, the right-hand side is the power the rail force must supply to hold
$\Ehat$ fixed: the control cost equals the failure of the Killing equation, i.e.\
the missing symmetry. This is why the non-stationary rail is legitimate but not
free---$\Ehat$ is controlled, not conserved. The optical reductions of $H$
reproduce the Randers/Zermelo structure
\cite{GibbonsRanders2009,CaponioWind}.

\paragraph{Geometric dictionary.} The indicatrix~\eqref{eq:ellipse} carries two
independent data---its \emph{center} $c=(c_r,c_\varphi)$ (the Zermelo wind) and
the \emph{eccentricity} of its semi-axes (the anisotropy of the local speed,
$a_r\neq a_\varphi$)---and each governs a distinct feature of the
brachistochrones, the same for every branch:
\begin{itemize}\itemsep2pt
\item eccentricity $\neq0$ (a spatial gradient $\partial_r f\neq0$)
$\Rightarrow$ the $t$ and $\tau$ branches split in depth, $n_t/n_\tau=E/f\neq1$;
\item azimuthal offset $c_\varphi\neq0$ (frame dragging)
$\Rightarrow$ rotational plunge inversion becomes possible;
\item radial offset $c_r\neq0$ (a non-stationary wind, $m'\neq0$)
$\Rightarrow$ teleological memory, $p_v\neq0$ along the path;
\item shrinking axes $R\to0$ (local speed $v\to0$) $\Rightarrow$ freezing and
capture.
\end{itemize}
Eccentricity and center offset are \emph{separate} levers: the first switches on
the $t/\tau$ split, the second the inversions. In particular the three branches
\emph{degenerate}---the $t$, $\tau$ and $\eta$ brachistochrones coincide---iff the
indicatrix is, up to a position-independent conformal scale, a rigid centered
circle ($c=0$, $a_r=a_\varphi$, radius depending on the flow parameter alone),
equivalently iff the optical cost ratios $n_t/n_\tau,\,n_t/n_\eta$ are
position-independent. This is exactly spatial homogeneity: FLRW is the unique
such case in our ladder (Sec.~\ref{sec:flrw}), whereas any gradient or wind
breaks it---even the static, non-rotating Schwarzschild ellipse, which is
centered ($c=0$) yet eccentric, already splits $t$ from $\tau$. The indicatrix
shape is thus the organizing thread of the paper: FLRW (circle) $\to$ Vaidya
(radial wind) $\to$ Thakurta--Kerr (angular wind) is a progression of ellipse
deformations, each switching on one further phenomenon.

\begin{figure}[tbp]
    \centering
    \includegraphics[width=\columnwidth]{Immagini/fig_indicatrici.png}
    \caption{Rail indicatrices (linear constraint $\cap$ normalization
    $=$ ellipse): the unifying formalism. FLRW gives shrinking circles (no
    wind, no eccentricity $\Rightarrow$ degenerate branches); Vaidya an ingoing
    radial wind ($f-E^2<0$, driven by $m(v)$ $\Rightarrow$ teleological memory);
    Thakurta--Kerr an angular wind (frame dragging $\Rightarrow$ rotational
    inversion), the ellipse shrinking with $A(\eta)$. Eccentricity switches on
    the $t/\tau$ split, the center offset the inversions.}
    \label{fig:indicatrici}
\end{figure}

\subsection{Selecting $W$: the Killing--CKV--Kodama hierarchy}
\label{sec:whierarchy}
The selector $W$ is fixed by a hierarchy of decreasing symmetry, each level
recovering the previous one in the appropriate limit.

\emph{(a) Killing.} If a timelike Killing vector $\xi$ exists, $W=\xi$ and the
right-hand side of Eq.~\eqref{eq:cost} vanishes identically: $-u\cdot\xi$ is
Noether-conserved and the rail is free. This is Perlick's stationary case
\cite{Perlick1991}.

\emph{(b) Conformal Killing.} If the spacetime is conformally stationary,
$g=\Omega^2\hat g$ with $\hat g$ stationary, the timelike conformal Killing
vector $\partial_\eta$ selects the conformal-time branch; $-u\cdot W$ is exactly
conserved for null rays and controlled up to the conformal factor for massive
rails. FLRW and constant-$A$ Thakurta--Kerr live at this level.

\emph{(c) Kodama.} In a spherically symmetric \emph{dynamical} spacetime with no
Killing or conformal Killing vector, the Kodama vector
$K^a=\varepsilon^{ab}\nabla_b r$ \cite{Kodama1980,Hayward1996} is the canonical
replacement: it is divergence-free (Eq.~\eqref{eq:kodama}), reduces to $\xi$ in
the static limit, and its flux $G^{ab}K_b$ is conserved with charge the
Misner--Sharp mass. Vaidya is the paradigm (Sec.~\ref{sec:kodama-energy}).

\emph{(d) Convention.} Absent even spherical symmetry the Kodama construction
fails and $W$ becomes a declared convention, with no distinguished normalization
of the cost one-form. Levels (a)--(c) are what make the non-stationary rail
canonical rather than arbitrary.


% ======================================================================
\section{FLRW: the degenerate base case}
\label{sec:flrw}
% ----------------------------------------------------------------------
For spatially flat FLRW, $g=a(\eta)^2\eta_{\rm Mink}$, the rail uses the
conformal Killing vector $\partial_\eta$; the local speed is
\begin{equation}
    v(\eta)=\sqrt{1-a(\eta)^2/\Ehat^{2}},
    \label{eq:flrwspeed}
\end{equation}
vanishing at the freezing barrier $a=\Ehat$.
\paragraph{Branch degeneracy.} Because $a=a(\eta)$ depends on conformal time
alone (spatial homogeneity), the optical index built from the conformal Killing
vector,
\begin{equation}
    n(\eta,r)=\frac1v=\frac{\Ehat}{\sqrt{\Ehat^2-a(\eta)^2}},
    \label{eq:flrwindex}
\end{equation}
is \emph{position-independent} at fixed $\eta$. The conformal-time functional
$T_\eta=\int n\,\dd\ell_{\rm flat}$ is then minimized by a flat straight line
(Fermat with a purely time-dependent index). The same curve minimizes $T_t$ and
$T_\tau$: with $\dd t=a\,\dd\eta$ and $\dd\tau=a\sqrt{1-\dot{\vec x}^2}\,\dd\eta$,
the monotone reparametrization $t(\eta)=\int a\,\dd\eta$ maps critical points to
critical points, and the $\tau$-cost differs from the $\eta$-cost only by the
same $\eta$-dependent weight. Hence
\begin{equation}
\begin{split}
    \arg\min T_t&=\arg\min T_\tau=\arg\min T_\eta\\
    &=\text{comoving straight line}.
\end{split}
    \label{eq:flrwdegen}
\end{equation}
The $t/\tau$ splitting that drives every later result is a gravitational-%
\emph{gradient} effect ($\partial_r f\neq0$); it is absent whenever the metric
functions depend on time only. FLRW is thus the degenerate base of the ladder:
it exhibits freezing at $a=\Ehat$ (Fig.~\ref{fig:flrw-kin})
but no branch structure (Fig.~\ref{fig:flrw-var}).

\begin{figure}[tbp]
    \centering
    \includegraphics[width=\columnwidth]{Immagini/fig_flrw_cinematica.png}
    \caption{de Sitter FLRW rail kinematics: local speed $v$ and rail redshift
    $E_{\rm phys}=\Ehat/a$; freezing at $a=\Ehat$.}
    \label{fig:flrw-kin}
\end{figure}


\begin{figure}[tbp]
    \centering
    \includegraphics[width=\columnwidth]{Immagini/fig_flrw_variazionale.png}
    \caption{Variational check: the comoving straight line minimizes
    \emph{both} $T_t$ and $T_\tau$ (branch degeneracy).}
    \label{fig:flrw-var}
\end{figure}


% ======================================================================
\section{Vaidya: dynamical black hole and the Kodama rail}
\label{sec:vaidya}
% ----------------------------------------------------------------------
The ingoing Vaidya metric $\dd s^2=-f\,\dd v^2+2\,\dd v\,\dd r+r^2\dd\Omega^2$,
$f=1-2m(v)/r$ \cite{Vaidya1951,LindquistSchwartzMisner1965}, has no timelike
Killing vector; its trapping structure is that of a dynamical horizon
\cite{AshtekarKrishnan2003,Booth2005,FaraoniBook}. The Kodama vector
\cite{Kodama1980,Hayward1996,AbreuVisser2010}
\begin{equation}
    K^a=\varepsilon^{ab}\nabla_b r = \partial_v ,\qquad
    \nabla_a K^a=0,
    \label{eq:kodama}
\end{equation}
is divergence-free without any Killing symmetry, and its charge is the
Misner--Sharp mass $m(v)$ \cite{MisnerSharp1964,HaywardMukohyama1999}.

\subsection{Kodama energy is the rail invariant}
\label{sec:kodama-energy}
From the rail constraint alone one finds the exact identity
$-g(T,T)=(f\,u^v-1)^2/\Ehat^2$, hence
\begin{equation}
    -\,u\cdot K \;=\; \Ehat \quad\text{identically along the brachistochrone.}
    \label{eq:kodama-conserved}
\end{equation}
This is \emph{controlled} conservation, not Noether: for a geodesic
$-u\cdot K$ would drift at rate $-m'(u^v)^2/r\propto m'$, exactly the
right-hand side of Eq.~\eqref{eq:cost} for $W=K$. Only the unnormalized
$K=\partial_v$ has zero control cost in the static limit; the normalized
$\hat K=K/\sqrt f$ does not (discriminating test), which is why the Kodama
vector, not a unit selector, is the canonical choice.

\paragraph{Indicatrix and Hamiltonians.} With $w\equiv\Ehat^2-f$ the constraint
$-u_v=\Ehat$ turns the admissible velocities at fixed $(v,r)$ into an ellipse,
\begin{equation}
    r'=(f-\Ehat^2)+\Ehat\sqrt w\,\cos\theta,\qquad
    \varphi'=\frac{\sqrt w}{r}\sin\theta,
    \label{eq:vaidya-ind}
\end{equation}
a Zermelo problem with ingoing radial wind $f-\Ehat^2<0$. Maximizing
$p\cdot\dot x$ over $\theta$ at fixed $p_\varphi=J$ gives the closed-form branch
Hamiltonians
\begin{align}
    H_v&=p_r(f-\Ehat^2)-1+\sqrt w\,\sqrt{\Ehat^2 p_r^2+J^2/r^2},
    \label{eq:Hv}\\
    H_\tau&=p_r(f-\Ehat^2)-\Ehat
        +\sqrt w\,\sqrt{(\Ehat p_r+1)^2+J^2/r^2}.
    \label{eq:Htau-vaidya}
\end{align}
The $\tau$-branch is the $v$-branch under $p_r\to p_r+1/\Ehat$ with unit cost
replaced by $\Ehat$. Non-stationarity is the explicit dependence
$\dd H/\dd v=\partial_v H\propto m'(v)$; the teleological momentum $p_v=-H$
satisfies $p_v(r_1)=0$ at the free arrival radius yet $p_v(r)\neq0$ along the
path, so the turning point \emph{anticipates future accretion}
(Fig.~\ref{fig:vaidya-traj}).

\begin{figure}[ht]
    \centering
    \includegraphics[width=\columnwidth]{Immagini/fig_kodama_conservazione.png}
    \caption{Kodama energy $-u\cdot K=\Ehat$ conserved along the
    $\tau$-brachistochrone to $6.7\times10^{-16}$ (top); the free-particle
    drift $\propto m'$ that the rail cancels (bottom).}
    \label{fig:kodama}
\end{figure}


\subsection{Phenomenology: penetration, timing, bounce}
\label{sec:vaidya-pheno}
Three features distinguish the dynamical hole from its frozen-mass
Schwarzschild limit.

\emph{Penetration.} The $\tau$-brachistochrone reaches the apparent horizon
$r=2m(v)$ only if $|J|$ lies below a threshold $J_c(v_0)$; accretion ($m'>0$)
opens a \emph{finite} window $J_c>0$, whereas evaporation ($m'<0$) closes it to
zero measure---an accretion/evaporation asymmetry with no static analogue. The
finite capture interval $J\in(0,\Jc]$ is the analogue of the Kerr $t$-branch
dichotomy, generated here by the \emph{dynamics}: grazing orbits linger near
periapsis while the growing horizon closes the gap.

\emph{Bounce.} The $r$-parametrization degenerates at periapsis
($\dd v/\dd r\to\infty$); reparametrizing by $v$ makes the Euler--Lagrange
problem regular at the turning point $r'=0$, so the trajectory is integrated
straight \emph{through} the bounce---the same extremal, via the exact change of
variable $(f v'-1)\dd r=(f-r')\dd v$. The $v$-flow reaches the periapsis exactly
where the $r$-flow diverges ($r_{\min}=3.0105$ at $v=7.29$ for $\mu\equiv m'=0.01$),
and the teleological memory $p_v(r)\propto m'$ measures the departure from the
frozen-mass Schwarzschild curve (Fig.~\ref{fig:vaidya-bounce}).

\emph{Timing and the orbit--horizon race.} Because $m$ grows along the path,
$r_{\min}$ depends on \emph{when} the particle arrives, the epoch $v_1$---a
genuinely non-stationary effect absent in Kerr. Strikingly, the absolute and
relative penetration have \emph{opposite} winners
(Fig.~\ref{fig:vaidya-timing}): under evaporation the absolute periapsis descends
but the horizon $2m(v)$ retreats faster, so the relative penetration
$r_{\min}/2m(v_{\rm peri})$ \emph{worsens} for later arrivals; under accretion the
periapsis rises yet the horizon grows faster, so late arrivals graze the relative
horizon. The orbit--horizon chase has reversed outcomes measured in $r$ versus in
horizon units.

\begin{figure}[ht]
    \centering
    \includegraphics[width=\columnwidth]{Immagini/fig_vaidya_traiettorie.png}
    \caption{$\tau$-brachistochrones Schwarzschild vs Vaidya, and $p_v(r)$:
    transversality $p_v(r_1)=0$, non-autonomous memory $\propto m'$.}
    \label{fig:vaidya-traj}
\end{figure}


\subsection{Plunge-radius law and the (absent) evaporative inversion}
\label{sec:vaidya-plunge}
The turning-point law has the static Schwarzschild form evaluated at the mass
at periapsis, $f_*=1-2m(v_{\rm peri})/r_{\min}$:
\begin{align}
    \tau:&\quad (\Ehat^2-f_*)J^2 = f_* r_{\min}^2 ,\\
    t:&\quad (\Ehat^2-f_*)J^2 = \Ehat^2 r_{\min}^2/f_* ,
    \label{eq:vaidya-turn}
\end{align}
plus a genuine $O(m')$ correction from the teleological momentum $p_v$. The
static ratio $J^2_t/J^2_\tau=\Ehat^2/f^2>1$ gives $r_{\min}^t<r_{\min}^\tau$
(``$t$ deeper''), robust for physical evaporation: \emph{there is no physical
evaporative inversion} (the apparent linear-model inversion is an artifact of
$m\le0$).

\begin{figure}[ht]
    \centering
    \includegraphics[width=\columnwidth]{Immagini/fig_vaidya_plunge_t_tau.png}
    \caption{Plunge radii $r_{\min}^{t,\tau}$ and gap vs $\mu=m'$; $t$ deeper,
    accretion widens the gap.}
    \label{fig:vaidya-plunge}
\end{figure}

\begin{figure}[ht]
    \centering
    \includegraphics[width=\columnwidth]{Immagini/fig_vaidya_no_inversione_evaporazione.png}
    \caption{No physical evaporative inversion: the linear-model sign change
    lies entirely in the unphysical $m<0$ region; with $m=e^{-\lambda v}>0$ the
    gap shrinks but never inverts.}
    \label{fig:vaidya-noinv}
\end{figure}

\subsection{Closed-form adiabatic correction and the horizon dilogarithm}
\label{sec:vaidya-adiabatic}
For slowly varying $m(v)$ the frozen-coefficient WKB scheme of
Sec.~\ref{sec:adiabatic} applies with $m(v)$ in place of the conformal factor and
$\dot m=\dd m/\dd v$ the small parameter. The frozen family is the Schwarzschild
brachistochrone---the $a\to0$ limit of the Thakurta--Kerr forms---on
$S=r(r-2m)DE\,[r^2(r-2m)-J^2 DE]$, $DE=(E^2-1)r+2m$, which keeps six distinct roots:
the constrained brachistochrone is genus \emph{two} even in Schwarzschild (only
geodesics are elliptic). The mass derivative reduces exactly (horizon factor
cancelling), $\partial_m F=N_m/S^{3/2}$,
$\partial_m\varphi_0=\mathcal{A}^m/\sqrt S+\sum_{k=0}^4 c_k^m U_k$, $U_k=\int
r^k/\sqrt S\,\dd r$, algebraic $c_k^m$ (verified to $10^{-15}$). The clock is the
advanced time: from $v=t+r_*$,
\begin{equation}
    \frac{\dd v}{\dd r}=\frac{E\,r^3}{\sqrt S}+\frac{r}{r-2m}
    \;\Rightarrow\; v=E\,U_3+r+2m\ln(r-2m),
    \label{eq:vaidya-clock}
\end{equation}
the tortoise term injecting a horizon logarithm. This is the advanced-time (native)
branch; the \emph{proper-time} branch uses the second-kind clock
$\dd\tau/\dd r=r^2(r-2m)/\sqrt S$ ($\eta=U_3-2m U_2$, $b=(0,0,-2m,1,0)$, no horizon
log) and closes by the same assembly~\eqref{eq:psi-split} with $Q_{kj}=c_k^m b_j-c_j^m
b_k$, all coefficients algebraic (verified to $10^{-14}$). The advanced-time
first-order correction integrates by parts into a fully explicit closed form (verified
to $9\times10^{-14}$),
\begin{equation}
    \frac{\delta\varphi}{\dot m}=A_m\,v-E\!\int\! A_m\frac{r^3}{\sqrt S}\dd r
    -\int\! A_m\frac{r}{r-2m}\dd r,\quad A_m\equiv\partial_m\varphi_0,
    \label{eq:vaidya-deltaphi}
\end{equation}
Writing $A_m=\mathcal{A}^m/\sqrt S-\alpha_0+\sum_k c_k^m U_k$ with the constant
$\alpha_0=\mathcal{A}^m(r_0)/\sqrt{S(r_0)}$, integrating by parts in each factor
gives the \emph{fully assembled} closed form (verified to $2\times10^{-13}$)
\begin{equation}
\begin{aligned}
    \frac{\delta\varphi}{\dot m}
    ={}&A_m\,v
    -E\Big[\mathcal{I}_{\rm poly}-\alpha_0 U_3
        +\tfrac12\sum_{k=0}^4 c_k^m\big(U_kU_3-W_{3k}\big)\Big]\\
    &{}-\Big[\mathcal{L}_{2m}-\alpha_0(v-E\,U_3)
        +\sum_{k=0}^4 c_k^m\big(r\,U_k-U_{k+1}+2m\,\mathcal{D}_k\big)\Big],
\end{aligned}
    \label{eq:vaidya-full}
\end{equation}
expanding into: the abelian integrals $U_k=\int r^k/\sqrt S\,\dd r$; the elementary
rational-log primitive $\mathcal{I}_{\rm poly}=\int\mathcal{A}^m r^3/S\,\dd r$; the
genus-two polylogarithms $W_{3k}=\int(U_3\dd U_k-U_k\dd U_3)$ at infinity (as in
Kerr); the \emph{second}-kind horizon integral
$\mathcal{L}_{2m}=\int\mathcal{A}^m r/[(r-2m)\sqrt S]\,\dd r$ (as
$\mathcal{A}^m(2m)\ne0$ the integrand is $\sim(r-2m)^{-3/2}$: a double pole with no
residue at the horizon Weierstrass point, closed by the Kleinian
$\zeta$/quasi-periods~\cite{BuchstaberEnolskii1997,BakerAbelian1897}, not Fay); the
genuinely third-kind object is the horizon letter $\dd r/(r-2m)$ (tortoise log, simple
pole residue two), closing as a Fay $\theta$-ratio~\cite{Fay1973} as $L$
[Eq.~\eqref{eq:thirdkind-theta}] but centred at $r=2m$; and
the \emph{horizon dilogarithms} $\mathcal{D}_k=\int U_k/(r-2m)\,\dd r$. Because $r=2m$
is a branch point of $S$, in $t=\sqrt{r-2m}$ one has $\dd r/(r-2m)=2\,\dd t/t$---a
simple pole of residue two, a third-kind $\dd\log$ letter \emph{at the horizon}---with
elementary primitive $\ln(r-2m)$, so one integration by parts closes $\mathcal{D}_k$
onto a single weight-two letter (verified to $10^{-12}$),
\begin{equation}
    \mathcal{D}_k=U_k\ln(r-2m)-\mathcal{G}_k,\qquad
    \mathcal{G}_k=\int\ln(r-2m)\,\frac{r^k}{\sqrt S}\,\dd r,
    \label{eq:horizon-dilog}
\end{equation}
$\mathcal{G}_k$ being the horizon-anchored genus-two dilogarithm---the exact
counterpart of the Kerr $W_{jk}$, same hyperelliptic-dilogarithm class (irreducible by
the theta-divisor obstruction), but sourced by the advanced-time divergence at the
horizon rather than a third-kind dipole at infinity. On the separatrix $|J|=\Jc$ the
sextic gains a double root at $r_d$, the genus drops to one
($\sqrt S=(r-r_d)\sqrt{Q_4}$), the $U_k$ reduce to elliptic integrals (machine
precision vs.\ direct genus-two quadrature), and $\mathcal{G}_k$ collapses to the
tabulated elliptic (Bloch--Wigner) dilogarithm on $E:w^2=Q_4$---the maximally
explicit reusable form. For $m=1,E=1.4$: $\Jc=7.0266$, $r_d=-3.3637$, branch points
$\{-2.083,0,2m,8.727\}$, $k^2=0.6067$, rectangular nome $\tau=0.90597\,i$ (Legendre
$iK'/K$ and Sage lattice agree to fifteen digits); the horizon $r=2m$ is a branch
point of $Q_4$, hence a two-torsion point carrying the elliptic $\dd\log$ letter. With
the Weierstrass $\sigma,\zeta$ of $E$ (base at the Weierstrass point, involution
$z\to-z$) the weight-one integrals close as third-kind elliptic integrals,
$U_0=\rho[\ln\sigma(z-z_d)-\ln\sigma(z+z_d)]+Cz$, residue $\rho=1/\sqrt{Q_4(r_d)}$ and
$C$ algebraic (derived, verified to $10^{-14}$); one further integration turns
$\mathcal{G}_k$ into the elliptic
dilogarithm~\cite{Bloch2000,BeilinsonLevin1994} $\int\ln\sigma\,\zeta\,\dd z$, with
single-valued projection the tabulated Zagier series~\cite{Zagier1990}
$D^E(z)=\sum_{n\ge0}D(q^n\zeta)-\sum_{n\ge1}D(q^n/\zeta)$. All coefficients
are algebraic; the transcendence lives in $W_{jk},\mathcal{G}_k$, so the correction is
reusable once $S$ and the $c_k^m$ are fixed.

\paragraph{Accretion--evaporation asymmetry.} At first order $\delta\varphi$ is linear
in $\dot m$, so $\dot m\to-\dot m$ na\"ively just flips its sign. The physical
asymmetry is subtler: accretion ($\dot m>0$) is \emph{ingoing} Vaidya, clocked by the
advanced time $v=t+r_*$; evaporation ($\dot m<0$) is \emph{outgoing} Vaidya, clocked by
the retarded time $u=t-r_*$. The frozen reduction $\partial_m F$ is identical (it
depends only on $S$); only the clock changes, and $v\leftrightarrow u$ flips the sign
of the tortoise. With $v=E\,U_3+r_*$, $r_*=(r-r_0)+2m\ln\frac{r-2m}{r_0-2m}$,
\begin{equation}
\begin{aligned}
    \frac{\delta\varphi}{\dot m}\Big|_{\rm accr}&=\mathcal{A}_\infty+\mathcal{B}_{\rm hor},
    \qquad
    \frac{\delta\varphi}{\dot m}\Big|_{\rm evap}=\mathcal{A}_\infty-\mathcal{B}_{\rm hor},\\
    \mathcal{A}_\infty&=\!\int\!\partial_m F\,E U_3\,\dd r,
    \qquad
    \mathcal{B}_{\rm hor}=\!\int\!\partial_m F\,r_*\,\dd r,
\end{aligned}
    \label{eq:vaidya-asymmetry}
\end{equation}
so the net asymmetry at equal $|\dot m|$ is $2|\dot m|\,\mathcal{B}_{\rm hor}$, carried
\emph{entirely by the tortoise}: the infinity-polylogarithm bulk $\mathcal{A}_\infty$
($E U_3$ sector) is common and cancels. The transcendental, direction-sensitive part of
$\mathcal{B}_{\rm hor}$ is the horizon term $2m\!\int\!\partial_m F\ln(r-2m)\dd r$,
which by parts is the horizon dilogarithm $\mathcal{D}_k$ of
Eq.~\eqref{eq:horizon-dilog}. Thus the accretion/evaporation asymmetry is the physical
face of the horizon weight-two piece---just as the ergosphere-penetration window is the
phenomenological face of the same horizon geometry. For $m=1,E=1.4,J=2.5$:
$\mathcal{A}_\infty=8.7276$, $\mathcal{B}_{\rm hor}=6.1127$ (horizon-log dilogarithm
sector $1.5301$), giving $\delta\varphi/\dot m|_{\rm accr}=14.840$ vs.\
$\delta\varphi/\dot m|_{\rm evap}=2.615$ (verified to $10^{-15}$).


% ======================================================================
\section{Thakurta--Kerr: conformal rotating compact object}
\label{sec:thakurta}
% ----------------------------------------------------------------------
Thakurta--Kerr is conformal Kerr, $g=A(\eta)^2 g_{\rm Kerr}$
\cite{Thakurta1981,Kerr1963}, one of a family of cosmological geometries
\cite{McVittie1933,SultanaDyer2005,Kaloper2010,NolanMcVittie1998,MelloMaciel2017,FaraoniBook}.
For constant $A$
the transfer theorem maps the problem to Kerr with $E_{\rm eff}=\Ehat/A$,
$J_{\rm eff}=J/A$; the horizon $\rp$ is rigid (conformally invariant), the
ergosphere $\re=2M$ is crossed regularly, while the freezing surface
$A^2 f=\Ehat^2$ breathes.

\paragraph{Indicatrix and Hamiltonians.} Equatorially, write
$\Delta=r^2-2Mr+s^2$, $\bar v^2=1-A^2f/\Ehat^2$, and
$\bar P=P+A^2(2Ms/r)^2/\Ehat^2$ with $P=r^2+s^2+2Ms^2/r$. The rail
$\Ehat=-u_\eta$ makes the indicatrix an ellipse with an \emph{azimuthal} wind
(frame dragging),
\begin{equation}
    \varphi'_0=\frac{2Ms}{r}\frac{\bar v^2}{\bar P},\qquad
    R^2=f\bar v^2+\bar P\varphi_0'^2=\frac{\Delta\,\bar v^2}{\bar P},
    \label{eq:tk-ind}
\end{equation}
the last equality using the Kerr identity $fP+(2Ms/r)^2=\Delta$. The
conformal-time, coordinate-time and proper-time Hamiltonians are
\begin{align}
    H_\eta&=p_\varphi\varphi'_0
        +R\sqrt{(\Delta/r^2)p_r^2+p_\varphi^2/\bar P}-1,
    \label{eq:Heta}\\
    H_t&=p_\varphi\varphi'_0
        +R\sqrt{(\Delta/r^2)p_r^2+p_\varphi^2/\bar P}-A,
    \label{eq:Ht-tk}\\
    H_\tau&=\tilde p_\varphi\varphi'_0
        +R\sqrt{(\Delta/r^2)p_r^2+\tilde p_\varphi^2/\bar P}-\frac{A^2 f}{\Ehat},
    \label{eq:Htau-tk}
\end{align}
with the conformal gravitomagnetic shift
$\tilde p_\varphi=p_\varphi-2MsA^2/(r\Ehat)$. The arrival-time branch $H_t$ differs
from $H_\eta$ only in the unit cost ($-1\to-A$) and reduces to it for $A=1$; it is
the branch compared with $\tau$ in the plunge inversion
(Sec.~\ref{sec:inversion}), where its turning potential
$H_{t0}=H_t|_{p_r=0}$ enters. All three are regular through the
ergosphere ($f=0\Rightarrow R^2=\bar v^2\Delta/\bar P>0$) and degenerate only at
the horizon $\Delta=0$, which is therefore conformally invariant---independent of
$A$ and $\Ehat$. The three branches share this ellipse: $t$ follows the same
curves as $\eta$ through $t=\int A\,\dd\eta$, while $\tau$ is Riemannian-pure, the
wind cancelling the gravitomagnetic term, so it survives the expansion
(Fig.~\ref{fig:th-branches}).

\paragraph{Non-autonomous optical reduction.} Solving the indicatrix for
$\dd\eta$ at fixed spatial displacement casts the $\eta$-branch as a closed-form
Randers metric, non-autonomous through the conformal factor,
\begin{align}
    \dd\eta &= n(\eta,r)\,\alpha_K+\beta_K,\qquad
    n=\frac{\Ehat}{\sqrt{\Ehat^2-A^2 f}},
    \label{eq:randers}\\
    \alpha_K^2 &= \frac{r^2\,\dd r^2}{f\Delta}+\frac{\Delta\,\dd\varphi^2}{f^2},
    \qquad
    \beta_K = -\frac{2Ms}{rf}\,\dd\varphi.
    \label{eq:randers-data}
\end{align}
The Riemannian data $\alpha_K$ and the one-form $\beta_K$ are \emph{exactly} those
of the \emph{null} Fermat (Randers) metric of Kerr \cite{GibbonsWerner2008}: all
of the mass, rail, and expansion physics is carried by the scalar index $n$,
which multiplies only the Riemannian part, while the frame-dragging one-form
$\beta_K$ is rigid---conformally invariant and independent of $\Ehat$. The
reduction is thus a Zermelo problem with a time-dependent landscape in which only
the metric term breathes, not the wind. It degenerates at the ergosphere ($f=0$,
where $1/F$ is intrinsic, as for stationary Kerr) even though the
Hamiltonian~\eqref{eq:Heta} remains regular there; the limits $\Ehat\to\infty$
(null Kerr optics), $A=1$ (stationary Perlick on Kerr), $s=0$
(Thakurta--Schwarzschild) and $M=s=0$ ($n=1/v$, FLRW) are all recovered.



\begin{figure}[ht]
    \centering
    \includegraphics[width=\columnwidth]{Immagini/fig_thakurta_kerr_rami.png}
    \caption{The three branches: $\tau$ bounces (mirror pair, no wind), $\eta$
    spirals (gravitomagnetic wind); breathing conformal scalars $n,k_\tau$.}
    \label{fig:th-branches}
\end{figure}

\paragraph{Branch optical metrics across the ladder.} The same Zermelo inversion
casts \emph{every} branch of \emph{every} metric in the paper as a Randers form
$\dd(\text{cost})=n\,\alpha+\beta$ (Table~\ref{tab:optical}), with the Riemannian
Fermat data $\alpha$ and wind one-form $\beta$ shared by the branches of a given
metric and only the scalar index $n$ set by the clock. Three structural facts
hold throughout and are verified to machine precision: (i) the proper-time branch
is always \emph{Riemannian-pure}, $\beta_\tau=0$---the wind (ingoing for Vaidya,
gravitomagnetic for Thakurta--Kerr) cancels in $\tau$; (ii) the arrival/native
branches carry the wind but share $\alpha$, and the arrival--proper split is the
universal ratio $n_t/n_\tau=E/f$ ($E=\Ehat$ for Vaidya, $E_{\rm eff}=\Ehat/A$ for
Thakurta--Kerr); (iii) non-stationarity makes only $\alpha,n$ \emph{breathe}
(through $m(v)$ or $A(\eta)$), while $\beta$ is rigid. FLRW is the degenerate top
of the ladder: $n$ is position-independent so all three branches share the same
(straight) geodesic. The $s=0$ limit of the Thakurta $\tau$-row reproduces the
static Vaidya $\tau$-row, and $M=s=0$ (or $m=0$) gives $n=1/v$, FLRW.

\begin{table*}[t]
    \centering
    \caption{Branch optical (Randers) metrics $\dd(\text{cost})=n\,\alpha+\beta$.
    Here $f=1-2m(v)/r$ (Vaidya) or $1-2M/r$ (Thakurta, with spin $s$,
    $\Delta=r^2-2Mr+s^2$), $w=\Ehat^2-f$, $w_A=\Ehat^2-A^2f$. The wind $\beta$
    acts on the native/arrival branch ($v$ for Vaidya, $\eta,t$ for Thakurta);
    $\tau$ is always Riemannian-pure ($\beta_\tau=0$). Native indices:
    $n_v=n_t^{\rm V}$ (Vaidya), $n_\eta=\Ehat/\sqrt{w_A}$ (Thakurta).}
    \label{tab:optical}
    \begin{ruledtabular}
    \begin{tabular}{lcccc}
    metric & $\alpha^2$ & $\beta$ & $n_\tau$ & $n_t$ \\
    \hline
    FLRW & $\dd r^2+r^2\dd\varphi^2$ & $0$
      & \multicolumn{2}{c}{$\Ehat/\sqrt{\Ehat^2-a^2}$ \ (degenerate)} \\[2pt]
    Vaidya & $\dd r^2+f\,r^2\dd\varphi^2$ & $\dd r/f$
      & $1/\sqrt{w}$ & $\Ehat/(f\sqrt{w})$ \\[2pt]
    Thakurta--Kerr
      & $\dfrac{r^2\dd r^2}{f\Delta}+\dfrac{\Delta\,\dd\varphi^2}{f^2}$
      & $-\dfrac{2Ms}{rf}\dd\varphi$
      & $\dfrac{A^2 f}{\sqrt{w_A}}$ & $\dfrac{A\Ehat}{\sqrt{w_A}}$ \\[3pt]
    \end{tabular}
    \end{ruledtabular}
\end{table*}

\subsection{Quasi-constants and drift}
\label{sec:quasiconst}
For constant $A$ the transfer theorem is exact at the level of conserved
quantities: $u_{\rm TK}=u_{\rm K}/A$ maps the $\Ehat$-rail onto the Kerr rail with
$E_{\rm eff}=\Ehat/A$, so every Kerr quasi-constant carries over by substitution,
\begin{equation}
    K^{\rm TK}=K^{\rm Kerr}(E_{\rm eff}),\qquad
    K_\tau=Q_{\rm std}+a^2 f_2\,p_\theta^2+O(a^4),
    \label{eq:Ktransfer}
\end{equation}
with $Q_{\rm std}$ the Carter constant and $f_2$ the next-to-leading polar
coefficient \cite{Carter1968,RosignoliKerr}. The polar libration that carries the
quasi-constant lives off the equatorial plane, governed by the full Hamiltonian
on a generic slice $\Sigma=r^2+s^2\cos^2\theta$,
\begin{equation}
    H_\tau=\tilde p_\varphi\varphi'_0
      +R\sqrt{\frac{\Delta}{\Sigma}p_r^2+\frac{p_\theta^2}{\Sigma}
      +\frac{\tilde p_\varphi^2}{\bar G}}-\frac{A^2 f_\Sigma}{\Ehat},
    \label{eq:H3d}
\end{equation}
with $f_\Sigma=1-2Mr/\Sigma$, $b=2Mar\sin^2\theta/\Sigma$,
$\bar G=G+A^2 b^2/\Ehat^2$, $\tilde p_\varphi=J-A^2 b/\Ehat$,
$\varphi'_0=b\bar v_\Sigma^2/\bar G$, $R^2=\bar v_\Sigma^2\Delta\sin^2\theta/\bar G$,
and the block identity $f_\Sigma G+b^2=\Delta\sin^2\theta$; it reduces to
Eqs.~\eqref{eq:Heta}--\eqref{eq:Htau-tk} on the equator. For $A=A(\eta)$ the
substitution~\eqref{eq:Ktransfer} is only the zeroth order: the transport equation
$\dd K/\dd\lambda=\{K,H\}+\partial_\eta K$ acquires an explicit source at order
$a^2\varepsilon$, $\varepsilon=A'/A$,
\begin{equation}
    \frac{\dd K}{\dd\eta}\bigg|_{\rm expl}=\varepsilon\,a^2\,S_1,\qquad
    S_1=2E_{\rm eff}^2\!\left[\cos^2\theta
        +\frac{M\cos2\theta\,p_\theta^2}{r\,DE_0^2}\right],
    \label{eq:S1}
\end{equation}
where $DE_0=r^2 w_{\rm eff}$. The counterterm $h$ solving $\{h,H_0\}=-S_1$
inherits the double pole at the \emph{freezing wall}
$r_w=2M/(1-E_{\rm eff}^2)$, so the dynamical quasi-constant is
\begin{equation}
    K^{\rm TK}(\eta)=K^{\rm Kerr}\!\big(E_{\rm eff}(\eta)\big)
        -a^2\!\int\!\varepsilon\,S_1\,\dd\eta,
    \label{eq:Kdyn}
\end{equation}
with closed toroidal averages
$\langle S_1\rangle=E_{\rm eff}^2(1-J^2/L^2)
+2ME_{\rm eff}^2|J|(L-|J|)/(r\,DE_0^2)$, $L^2=p_\theta^2+J^2/\sin^2\theta$. In
Kerr ($E>1$) the wall is absent; in Thakurta--Kerr with $E_{\rm eff}<1$ it enters
from the far field and advances inward as $A$ grows, becoming a \emph{global
attractor}: the bound $\tau$-libration ceases to exist, so the dynamical
quasi-constants live on scattering arcs ($E_{\rm eff}>1$) or on the $t$-branch.

\subsection{Semi-analytic first-order adiabatic orbit shape}
\label{sec:adiabatic}
The transport equation~\eqref{eq:Kdyn} controls the quasi-constant; the
\emph{orbit shape} $\varphi(r,A)$ itself admits a companion first-order adiabatic
form~\cite{BornFock1928,BenderOrszag1999} that is explicit up to a single higher
transcendental. Writing the frozen integrand as
$F(r;E)=\dd\varphi/\dd r$ (Eq.~\eqref{eq:t-K} for the $t$-branch,
Eq.~\eqref{eq:tau-sextic} for $\tau$) and linearizing $E_{\rm eff}=\Ehat/A$ in
$A(\eta)=A_0(1+\varepsilon\eta+\dots)$, $\varepsilon=A'/A$, gives
\begin{equation}
    \varphi(r,A)=\underbrace{\varphi_0(r;E_{\rm eff})}_{\text{frozen, Kleinian}}
    +\varepsilon\big[\mathcal{C}(r)+\psi(r)\big]+O(\varepsilon^2),
    \label{eq:phi-adiab}
\end{equation}
where the first-order coefficient is the elapsed-clock--weighted energy
derivative $\delta\varphi=-\Ehat\!\int\!\partial_E F\,\eta\,\dd r$, and $\eta(r)$
is the clock the orbit accumulates along the frozen slice: coordinate time
$\eta=t(r)$ on the $t$-branch, proper time $\eta=\tau(r)=\int\!L_\tau\,\dd r$ with
$L_\tau=\sqrt{Q/w}$, $Q=\Delta\varphi'^2-2a\beta\varphi'+r^2/(r^2+a^2)$, on the
$\tau$-branch (the extra $A^{-2}$ of the proper-time functional affects the
\emph{timing}, not the shape, so Eq.~\eqref{eq:phi-adiab} is common to both
branches).

The energy derivative reduces to the same depth-one class as $\varphi_0$ itself, by
the classical reduction of abelian differentials to first/second
kind~\cite{BakerAbelian1897,GriffithsHarris1978}.
Solving $\partial_E(F)=\dd_r[\mathcal{A}(r)/\mathcal{R}]
+\sum_{k=0}^{4}c_k\,r^k/\mathcal{R}$ for a degree-six numerator $\mathcal{A}(r)$
and constants $c_k$ ($\mathcal{R}$ the branch sextic radical, $\sqrt{R_6}$ or
$\sqrt{S}$) yields
\begin{equation}
    \partial_E\varphi_0=\frac{\mathcal{A}(r)}{\mathcal{R}}
    +\sum_{k=0}^{4}c_k\!\int\!\frac{r^k\,\dd r}{\mathcal{R}},
    \label{eq:dEphi0}
\end{equation}
with \emph{no} third-kind (horizon-pole) contribution: the second-kind reduction
closes with vanishing polar residues for both branches (verified to machine
precision). Hence $\mathcal{C}(r)=-\tfrac12\Ehat\,\partial_E\varphi_0\,\eta(r)$ is
a product of depth-one objects, closed in the Kleinian/abelian class.

\emph{Clock differentials.} The clock $\eta$ enters $\mathcal{C}$ through its
differential, closed on the branch curve. On the $\tau$-branch the proper-time
density reduces to a rational second-kind differential,
\begin{equation}
    \frac{\dd\tau}{\dd r}=\frac{r^2(r-2M)}{\sqrt{S}},
    \label{eq:clock-tau}
\end{equation}
the frame-dragging term cancelling identically ($10^{-16}$). The $t$-branch
coordinate-time density is \emph{also entirely on-curve}: the identity
$(r-2M)Q_2+DE(J(r-2M)+2Ma)^2=E^2r^3\Delta$ (with $E^2-f=DE/r$) collapses the Randers
wind, $n_t\alpha=E^2r^3/(f\sqrt{R_6})$, and
\begin{equation}
    \frac{\dd t}{\dd r}=\frac{\rho_t(r)}{\sqrt{R_6}},\quad
    \rho_t=\frac{E^2r^3-(2Ma/r)\mathcal K}{f}=P_3(r)+\frac{R_\Delta(r)}{\Delta},
    \label{eq:clock-t}
\end{equation}
the ergosphere pole $r=2M$ removable, third-kind poles only at the horizons. The
second-kind part $P_3=E^2r^3+2E^2Mr^2+4E^2M^2r+8E^2M^3-2(E^2-1)JMa$ gives the clock
vector $b=(8E^2M^3-2(E^2-1)JMa,4E^2M^2,2E^2M,E^2,0)$ ($t$-analogue of the $\tau$-branch
$b=(0,0,-2M,1,0)$), and $R_\Delta$ is the degree-one horizon numerator (explicit in
\texttt{reproduce\_reductions.py}). \emph{Neither} branch leaves the curve (verified
against the Hamiltonian flow to $10^{-10}$); the $t$-branch merely adds a third-kind
horizon letter---the apparent Doran double cover of a first analysis was spurious.
Table~\ref{tab:branch-adiab} collects the branch data: the two solutions share the
structure~\eqref{eq:phi-adiab} and differ only in the branch curve and the clock.
\begin{table*}[t]
    \centering
    \caption{The two branches of the adiabatic solution~\eqref{eq:phi-adiab}. They
    share $\varphi=\varphi_0+\varepsilon[\mathcal{C}+\psi]$;
    only the curve $\mathcal{R}^2$, the frozen numerator $K$ ($F=K/\mathcal{R}$),
    and the clock $\dd\eta/\dd r$ change. $\Delta=r^2-2Mr+a^2$, $DE=(E^2-1)r+2M$,
    all at $E=E_{\rm eff}=\Ehat/A$; $\varphi_0,\partial_E\varphi_0$ are
    Kleinian/abelian on $y^2=\mathcal{R}^2$, and $\psi=\tfrac12\Ehat\sum Q_{kj}W_{kj}$
    has algebraic coefficients times genus-two polylogarithms
    [Eq.~\eqref{eq:psi-split}].}
    \label{tab:branch-adiab}
    \begin{ruledtabular}
    \begin{tabular}{lll}
         & $t$-branch (coord.) & $\tau$-branch (proper)\\
        \hline
        curve $\mathcal{R}^2$ & $R_6=r\,Q_2\,DE$ & $S=r(r{-}2M)DE[r\Delta{-}J^2DE]$\\
        numerator $K$ & $r\,DE(J(r{-}2M){+}2Ma)/\Delta$ & $J\,r(r{-}2M)DE/\Delta$\\
        clock $\dd\eta/\dd r$ & $\rho_t/\sqrt{R_6}$ (2nd$+$3rd)
            & $r^2(r-2M)/\sqrt{S}$ (2nd)\\
        off-curve cover & none & none\\
        timing prefactor & $1$ & $A^{-2}$\\
    \end{tabular}
    \end{ruledtabular}
\end{table*}

\emph{The correction $\psi$.} By parts and the shuffle identity~\cite{Chen1977}
$\rho+\tilde\rho=\partial_E\varphi_0\,\eta$
($\rho=\int\partial_E\varphi_0\,\dd\eta$,
$\tilde\rho=\int\eta\,\partial_E F\,\dd r$) is closed, and the antisymmetric
remainder
\begin{equation}
    \psi(r)=\tfrac12\Ehat\,(\rho-\tilde\rho)
    \label{eq:psi-irr}
\end{equation}
is a depth-two iterated integral. Expand $\partial_E\varphi_0$ and $\eta$ in the
cohomology basis $U_k=\int r^k/\mathcal{R}\,\dd r$: the reduction
$\partial_E F=N/\mathcal{R}^3$, $N=EJ\,r^4(r-2M)^2 DE$ (the horizon factor $\Delta$
cancels) gives $\partial_E\varphi_0=\mathcal{A}/\mathcal{R}+\sum_{k=0}^4 c_k U_k$,
and on $\tau$ $\eta=U_3-2M\,U_2$, i.e.\ $b=(0,0,-2M,1,0)$. The complete first-order
coefficient $\delta\varphi=\mathcal{C}+\psi=-\Ehat\int\partial_E F\,\eta\,\dd r$
reorganizes, with no term discarded, into
\begin{equation}
    \delta\varphi=
    \underbrace{\tfrac12\Ehat\!\sum_{k<j}Q_{kj}\,W_{kj}}_{\text{weight two}}
    -\underbrace{\Ehat\,\eta\big(\tfrac{\mathcal{A}}{\mathcal{R}}
    +\tfrac12\!\sum_k c_k U_k\big)}_{\text{clock}\times\text{abelian}}
    +\underbrace{\Ehat\!\sum_j b_j\,\mathcal{I}_j}_{\text{elementary}},
    \label{eq:psi-split}
\end{equation}
$Q_{kj}=c_k b_j-c_j b_k$, $W_{kj}=\int(U_k\dd U_j-U_j\dd U_k)$ the genus-two
polylogarithms, $\mathcal{I}_j=\int\mathcal{A}\,r^j/\mathcal{R}^2\,\dd r$ elementary
(rational, evaluating to logarithms). Every piece is explicit---the second group a
product of the clock with the reduced abelian integrals, the third an elementary
logarithmic remainder; none is spurious. The $Q_{kj}$ are \emph{algebraic}
(rational in $M,a,E,J$ through the $c_k$); we verified~\eqref{eq:psi-split} to
$4\times10^{-14}$ with all coefficients given, not fitted, and $Q_{01}=0$ (the clock
has no first-kind part). On the $t$-branch the same holds with
$\mathcal{R}^2=R_6$: the clock is likewise on-curve, its second-kind vector $b$
[Eq.~\eqref{eq:clock-t}] feeding the $W_{kj}$ of $R_6$ and its horizon part
$R_\Delta/\Delta$ the third-kind $\ln\sigma$-quotients at $r_\pm$---no double cover.
The full $t$-branch assembly reproduces the direct
$-\Ehat\!\int\!\partial_E F\,t\,\dd r$ to $9\times10^{-13}$.
The $W_{kj}$ are the genus-two hyperelliptic
polylogarithms of $y^2=\mathcal{R}^2$, weight-two iterated
integrals~\cite{Chen1977} in the (elliptic/higher-genus) polylogarithm class of
Beilinson--Levin and Brown--Levin~\cite{BeilinsonLevin1994,BrownLevin2011}: they do
not reduce to a weight-one Kleinian $\zeta,\sigma$
form~\cite{BakerAbelian1897,BuchstaberEnolskii1997}, since for genus two the Abel
image of single points lies on the theta divisor $\Theta=W_{g-1}$~\cite{MumfordTata}
(where $\zeta,\wp$ are singular). The single third-kind piece---the clock's dipole at
infinity, residue $\rho_0=M/(E^2-1)^{3/2}$---integrates to
\begin{equation}
    L(u)=\log\frac{\theta[\delta](u-e_+)}{\theta[\delta](u-e_-)}
    \label{eq:thirdkind-theta}
\end{equation}
(Fay prime form~\cite{Fay1973}, odd $\delta$, $e_\pm$ the Abel images of infinity;
$L(u(r))$ differentiated back reproduces the third-kind differential to
$5\times10^{-8}$). Like $\mathrm{Li}_2$, the $W_{kj}$ are their own closed form
(Kronecker--Eisenstein series~\cite{BeilinsonLevin1994}); the weight-two content is genuine and irreducible, the coefficients
algebraic.

Table~\ref{tab:adiab-valid} compares the hybrid
form~\eqref{eq:phi-adiab}--\eqref{eq:psi-irr} (closed $\mathcal{C}$ plus numerical
$\psi$) against the directly integrated non-autonomous coefficient
$-\Ehat\!\int\partial_E F\,\eta\,\dd r$, for both branches ($M=1$, $a=0.9$,
$\Ehat=1.4$; $t$-branch $J=6$, $\tau$-branch scattering $J=2.5$). The agreement is
at the discretization floor and scales linearly with $\varepsilon$, confirming
that the split is an exact reorganization (integration by parts), not an
approximation. Figure~\ref{fig:phi-closed} shows the same comparison as orbit
shapes; the frozen family itself penetrates the ergosphere for suitable angular
momenta, turning just outside the horizon and being expelled at an adiabatic
threshold (Fig.~\ref{fig:phi-penetrating}).
\begin{table}[htbp]
    \centering
    \caption{Hybrid semi-analytic adiabatic form versus full non-autonomous
    first-order coefficient, $\max_r|\varphi_{\rm hybrid}-\varphi_{\rm full}|$.
    Both branches; the residual is the numerical-quadrature floor.}
    \label{tab:adiab-valid}
    \begin{tabular}{lcc}
        \hline\hline
        $\varepsilon=A'/A$ & $t$-branch ($\eta=t$) & $\tau$-branch ($\eta=\tau$)\\
        \hline
        $0.005$ & $1.4\times10^{-6}$ & $2.5\times10^{-8}$\\
        $0.02$  & $5.6\times10^{-6}$ & $1.0\times10^{-7}$\\
        $0.06$  & $1.7\times10^{-5}$ & $3.0\times10^{-7}$\\
        \hline\hline
    \end{tabular}
\end{table}

\begin{figure*}[tbp]
    \centering
    \includegraphics[width=\textwidth]{Immagini/fig_phi_closed_validation.pdf}
    \caption{Closed-form adiabatic orbit shape $\varphi(r,A)$ against direct
    integration, both branches. (a,b): the closed form
    $\varphi_0+\varepsilon\,\delta\varphi$ (dashed), $\delta\varphi$ from
    Eq.~\eqref{eq:psi-split}, overlies the numerically integrated orbit (solid) over
    the whole arc through the turning point (marker), winding near $r_e$. (c): the
    residual scales linearly with $\varepsilon=A'/A$ ($t$-branch $J=6$, $\tau$-branch
    $J=2.5$). All coefficients are algebraic and given, not fitted.}
    \label{fig:phi-closed}
\end{figure*}

\begin{figure*}[tbp]
    \centering
    \includegraphics[width=\textwidth]{Immagini/fig_phi_penetrating.pdf}
    \caption{Ergosphere-penetrating $t$-branch orbits ($J=3<J_c^+$), frozen family.
    (a): as $A$ grows the brachistochrone enters the ergosphere ($r<r_e$),
    co-rotates, and turns ($\circ$) in $r_+<r_{\rm t}<r_e$---outside the horizon---%
    before exiting; the rail turning $R_6=0$ keeps it above $r_+$. (b): the
    penetration depth $r_{\rm t}(A)$ stays near $r_+$ until a threshold
    $A_c\!\simeq\!1.5$, beyond which the orbit no longer penetrates.}
    \label{fig:phi-penetrating}
\end{figure*}

\emph{Penetration phase diagram.} The frozen ergosphere dynamics is set by two
analytic thresholds (Fig.~\ref{fig:penetration-phase}). For $E_{\rm eff}<1$ a
freezing wall $r_w=2M/(1-E_{\rm eff}^2)$ advances inward; the launch radius $r_0$ is
swallowed at $A_c^{\rm wall}=\Ehat/\sqrt{1-2M/r_0}$, beyond which no orbit issues
from $r_0$. Below the wall the $t$-orbit penetrates the ergosphere iff its outer
turning lies inside $r_e$; the prograde boundary $J_c^+(A)=2M^2/a+a+aA^2/(2\Ehat^2)$
matches the numerics to $10^{-15}$. The window $(J_c^-,J_c^+)$ splits into a plunge
core and a thin penetrate-and-bounce rim below $J_c^+$, where the orbit dips inside
$r_e$, co-rotates, and turns in $r_+<r_{\rm t}<r_e$ before exiting
(Fig.~\ref{fig:bounce}); the breathing factor expels orbits from the ergosphere at
$A_c^{\rm wall}$, a purely non-stationary effect.

\begin{figure}[tbp]
    \centering
    \includegraphics[width=\columnwidth]{Immagini/fig_penetration_threshold.pdf}
    \caption{Ergosphere-penetration phase diagram of the frozen $t$-branch in
    $(A,J)$ ($M=1$, $a=0.9$, $\Ehat=1.4$, $r_0=12$): plunge, penetrate-and-bounce,
    scattering, forbidden (beyond $A_c^{\rm wall}=1.534$, dashed). $J_c^+(A)$ is
    analytic (solid, verified to $10^{-15}$).}
    \label{fig:penetration-phase}
\end{figure}

\begin{figure}[tbp]
    \centering
    \includegraphics[width=\columnwidth]{Immagini/fig_penetrate_bounce_orbit.pdf}
    \caption{Penetrate-and-bounce $t$-brachistochrone ($J=3.2$): enters the
    ergosphere, co-rotates (red), turns at $r_{\rm t}=1.87$ (inside $r_e$, outside
    $r_+$), and exits; $R_6=0$ keeps it above $r_+$.}
    \label{fig:bounce}
\end{figure}




% ======================================================================
\section{Equatorial closed forms and the ergosphere}
\label{sec:closed}
% ----------------------------------------------------------------------

All closed forms in this section and in App.~\ref{app:explicit} are obtained at
\emph{frozen conformal factor}: constant $A$, equivalently stationary Kerr with
the substitution $E\to E_{\rm eff}=\Ehat/A$ (and $a\to s$, $J\to J/A$) in every
parameter (Sec.~\ref{sec:thakurta}). They are therefore the exact
\emph{instantaneous} orbit shapes on a Thakurta--Kerr slice of given $A$; the
running $A(\eta)$ makes the full problem non-autonomous---the quasi-constant
drift of Sec.~\ref{sec:quasiconst}, $K^{\rm TK}(\eta)=K^{\rm Kerr}(E_{\rm
eff}(\eta))+O(A'/A)$---so a genuinely time-dependent trajectory is recovered by
letting $E_{\rm eff}=E_{\rm eff}(\eta)$ run through this family, not by a single
autonomous curve.

\subsection{Separatrix in Weierstrass functions}
\label{sec:separatrix}
From the doranTau factorization $\Delta-\Jc^2 w=f(r^2+c^2)$, $c=a/E$
\cite{RosignoliKerr}, the
equatorial $\tau$-separatrix $\varphi(r)$ is genus~1, closed in the
Weierstrass functions $\wp,\sigma,\zeta$
\cite{WhittakerWatson,DLMF,Hackmann2008,Hackmann2010}: the angle is linear in the
uniformizing variable $z(r)=\wp^{-1}(A/r+B)$ plus two $\sigma$-ratios
[App.~\ref{app:explicit}, Eq.~\eqref{eq:sep-phi}], regular \emph{through} the
ergosphere because the factor $r^2+c^2$ divides the turning quartic and removes
$\re$ from the radicand. The Boyer--Lindquist form winds logarithmically at the
horizon; the Doran shift---a second, quasi-lemniscatic elliptic
curve---cancels it, so the ergosphere-penetrating separatrix
$\varphi_D=\varphi_{\rm BL}+\varphi_{\rm shift}$ is closed on two elliptic curves.

What controls the reduction is the \emph{genus} of the turning radical, not the
rationality of $\mathcal K$. Being Riemannian-pure
(Sec.~\ref{sec:thakurta}, the gravitomagnetic wind cancels), the $\tau$-branch
factorizes as doranTau---the factor $r^2+c^2$ dividing the quartic---dropping a
generically genus-2 reduction to genus~1. Generic $t$-branch orbits carry the
wind and have the sextic radical
$\mathcal R(r)=r\,Q_2(r)\,[(E^2-1)r+2M]$ (genus~2, hyperelliptic), with no such
cancellation. The $t$-brachistochrone nonetheless has its own genus-1 closed-form
separatrix, by a \emph{different} mechanism: at the marginal-capture (saddle-node)
momentum $J_c^-$ [Eq.~\eqref{eq:Jc-minus}] the sextic develops a \emph{double
root}, $\mathcal R=(r-r_*)^2 Q_4$ with $Q_4$ quartic (whereas the prograde
ergosphere crossing $J_+^t$ has only a simple root and stays genus~2). The double
root drops the radical to $Q_4$, and $\varphi(r)$ closes in the same
Weierstrass $\wp,\sigma,\zeta$ form as the $\tau$-separatrix
(App.~\ref{app:explicit}). Both separatrices are thus genus~1---the $\tau$ one
via the doranTau factorization, the $t$ one via a double root---while generic
(non-separatrix) $t$-orbits are genuinely genus~2, closed only in genus-2
Kleinian functions.
The Kerr geodesic structure and its constants are classical
\cite{Carter1968,Chandrasekhar}. Validated end-to-end to
$2.8\times10^{-10}$ through the ergosphere, and general across parameters.

\begin{figure}[ht]
    \centering
    \includegraphics[width=\columnwidth]{Immagini/fig_separatrix_3traiettorie.png}
    \caption{Prograde $\tau$-branch separatrix $J=+\Jc$ (the single penetrating
    value): ODE, quadrature and evaluated Weierstrass agree to
    $2.8\times10^{-10}$; the PMP-$\eta$ terminates at $\re$.}
    \label{fig:sep3}
\end{figure}


\subsection{Trichotomy and the ergosphere cusp}
\label{sec:trichotomy}
For the equatorial $\tau$-branch, $\dd\varphi/\dd r\sim K\sqrt{r-\re}$ as
$r\to\re$ (because $\sqrt f\to0$ while the turning radicand stays finite),
giving
\begin{equation}
    \varphi-\varphi_e \sim \tfrac{2}{3}K\,(r-\re)^{3/2},
    \label{eq:cusp}
\end{equation}
a cusp with radial tangent. The trichotomy in $|J|$ (sign of $a^2-J^2E^2$):
smooth periapsis above $\re$ ($|J|>\Jc$); marginal contact with $\re$
($|J|=\Jc=a/E$, separatrix); cusp reflection at $\re$ ($|J|<\Jc$). Only the
\emph{prograde} marginal $J=+\Jc$ actually crosses: inside $\re$ frame dragging
forces co-rotation ($\dd\varphi/\dd\lambda>0$), so the retrograde marginal
$J=-\Jc$, which would counter-rotate, reflects with a cusp instead. The crossing
(prograde) branch is continued through $\re$ to the horizon in Doran
(horizon-penetrating) coordinates
\cite{Doran2000,Natario2009,MartelPoisson2001},
$\varphi_D=\varphi_{\rm BL}-\varphi_{\rm shift}$.

\paragraph{Retrograde branch and the $\tau/t$ asymmetry.} The trichotomy above is
a $\tau$-branch statement. The set that \emph{reaches} $\re$ is symmetric in the
sign of $J$---the doranTau factorization depends on $J$ only through $J^2$, so the
ingoing orbit reaches $\re$ for $|J|\le\Jc=a/E$, prograde and retrograde sharing
the same threshold despite frame dragging---but the set that \emph{penetrates} is
not: only the single prograde $+\Jc$ crosses, the co-rotation constraint breaking
the $\pm J$ symmetry at the wall. The $t$-branch is not protected. Its turning function (from
$H_\eta=0$ at $p_r=0$, with $b=2Ma/r$) is
\begin{equation}
    N(r)=\big(\bar P-Jb\bar v^2\big)^2-J^2\Delta\bar v^2,
    \label{eq:tbranch-turn}
\end{equation}
which is only linear-plus-quadratic in $J$. At the ergosphere
$N(\re)=\bar P_e\,(\bar P_e-2aJ)$ with
$\bar P_e=\bar P(\re)=4M^2+2a^2+a^2/E^2$, so the $t$-orbit crosses $\re$ at the
single \emph{prograde} value
\begin{equation}
    J_+^{t}=\frac{\bar P_e}{2a}=\frac{4M^2+2a^2+a^2/E^2}{2a};
    \label{eq:Jplus-t}
\end{equation}
there is no retrograde $r_e$-crossing. Instead the retrograde $t$-capture is
bounded by the \emph{marginal} (double-root) condition $N=N'=0$ at a radius
$r_*>\re$---the massive analogue of the retrograde photon sphere---which closes
to
\begin{equation}
    J_c^{-}=-\,\frac{\bar P}{\bar v\,(\sqrt\Delta-b\bar v)}\bigg|_{r_*}<0 .
    \label{eq:Jc-minus}
\end{equation}
For $a=0.9$, $E=1.2$ (and $A=1$) these give $J_+^{t}=3.43$ and $J_c^{-}=-8.05$
($r_*=3.51$): the $t$-branch captures the wide, asymmetric band
$J\in[J_c^{-},J_+^{t}]$, far more capture-prone on the retrograde side. A
moderately retrograde orbit ($J=-5$, say) is therefore captured on the
$t$-branch but scattered on the $\tau$-branch (Fig.~\ref{fig:master-pen}). For
constant $A$ the results carry
over by $E\to\Ehat/A$, $a\to s$, $J\to J/A$; in particular
$\Jc^{\tau}=sA^2/\Ehat$.

\begin{figure}[t]
    \centering
    \includegraphics[width=\columnwidth]{Immagini/fig_master_penetration_taut.png}
    \caption{Ergosphere behaviour of the two branches versus $J$
    (Table~\ref{tab:penetration}). The $\tau$-branch (bottom bar) is a
    \emph{trichotomy}: smooth bounce for $|J|>J_c$, cusp bounce at $r_e$ for
    $|J|<J_c$ (red), and a \emph{single} penetrating value $J=+J_c$ (star). The
    $t$-branch (top bar) is a \emph{dichotomy}: it penetrates the whole interval
    $(J_c^-,J_+^t)$ (green), reflecting smoothly outside.}
    \label{fig:master-pen}
\end{figure}

\begin{figure}[ht]
    \centering
    \includegraphics[width=\columnwidth]{Immagini/fig_thakurta_cuspide_ergosfera.png}
    \caption{Behaviour at the ergosphere: cusp reflection ($|J|<\Jc$), smooth
    periapsis ($|J|>\Jc$), and the \emph{prograde} marginal $J=+\Jc$ crossing to
    $\rp$ in Doran coordinates---only $+\Jc$ crosses, the retrograde $-\Jc$
    cusping instead. Power law $2/3$ (cusp) vs $2$ (periapsis).}
    \label{fig:th-cusp}
\end{figure}

\subsection{Plunge inversion: rotational and conformal}
\label{sec:inversion}
Two mechanisms invert the $t/\tau$ ordering. Rotationally, in the $(a,J)$ plane,
$\rho(r_*)^2=H_{t0}/H_{\tau0}$, extending the stationary Kerr analysis
\cite{RosignoliKerr}. Conformally, the inversion locus is
\begin{equation}
    \mathrm{cond}(r,A)=A^2 b\,Q-\bar P(\Ehat-A^2 f)=0,\quad
    Q=b\bar v^2+\sqrt{\bar v^2\Delta},
    \label{eq:condinv}
\end{equation}
with $\mathrm{cond}(1)<0$ and
$\mathrm{cond}(\Af)=\Ehat r\Delta(\Ehat-1)/(r-2M)>0$, so by the intermediate
value theorem $A_{\rm inv}\in(1,\Af)$: the conformal inversion always exists
and is reachable (far field $A_{\rm inv}\to\sqrt{\Ehat}$).

\begin{figure}[ht]
    \centering
    \includegraphics[width=\columnwidth]{Immagini/fig_thakurta_kerr_plunge_t_tau.png}
    \caption{Rotational inversion in the $(a,J)$ plane; ODE locus vs analytic
    $\rho(r_*)^2=H_{t0}/H_{\tau0}$.}
    \label{fig:th-plungeAJ}
\end{figure}

\begin{figure}[tbp]
    \centering
    \includegraphics[width=\columnwidth]{Immagini/fig_thakurta_kerr_inversione_AJ.png}
    \caption{Conformal inversion in the $(J,A)$ plane: $\Delta r$ changes sign;
    analytic $\mathrm{cond}=0$ matches the numerical contour; reachable
    ($A_{\rm inv}<\Af$).}
    \label{fig:th-invAJ}
\end{figure}


\subsection{Fixed-endpoint comparison and protocol dependence}
\label{sec:bvp}
The statement ``branch $X$ plunges deeper'' is protocol dependent. With
\emph{fixed initial data} (same $r_0,J,p_r$) the $t$-branch is deeper
($r_{\min}^t<r_{\min}^\tau$). With \emph{fixed endpoints} (the genuine
two-point brachistochrone), the two branches join the same $A,B$ but are
distinct paths, and the ordering reverses: $r_{\min}^\tau<r_{\min}^t$. The
optical-index ratio $n_t/n_\tau=E/f>1$ penalizes deep radii more heavily for
the $t$-branch, so the $t$-brachistochrone bows outward while the $\tau$-branch
dips. This must be specified when comparing depths.

\begin{figure}[ht]
    \centering
    \includegraphics[width=\columnwidth]{Immagini/fig_bvp_estremi_fissi.png}
    \caption{Fixed-endpoint $t$ and $\tau$ brachistochrones (Schwarzschild,
    $E=1.2$): symmetric (top) and asymmetric (bottom) endpoints. Both branches
    join the same $A,B$; the $\tau$-branch dips deeper.}
    \label{fig:bvp}
\end{figure}

Crucially, neither driver of same-launch inversion survives the
fixed-endpoint constraint. Scanning the spin $a$ over the entire
Kerr flow with fixed symmetric endpoints $(r_0,\pm\Phi)$ leaves
$\Delta r=r_{\min}^t-r_{\min}^\tau$ strictly positive
($\Delta r\in[+0.45,+2.06]$ on $a\in[0.05,0.95]$); the map is nearly flat
in $a$, its gradient set almost entirely by the endpoint angle
$\Phi$ (encoded by $J$). The conformal scan (via the transfer
$E_{\rm eff}=\hat E/A$) is likewise everywhere positive
($\Delta r\in[+0.002,+3.711]$), as forced by the theorem $n_t/n_\tau=E/f>1$.
Table~\ref{tab:protocol} summarizes the resulting protocol dependence:
inversion is a property of the \emph{comparison protocol}, not an invariant
of the two brachistochrones---it requires the freedom to let the endpoint
move.

\begin{figure}[tbp]
    \centering
    \includegraphics[width=\columnwidth]{Immagini/fig_colormap_spin_estremi_fissi.png}
    \caption{Spin scan at fixed symmetric endpoints (Kerr, Hamiltonian
    flow): $\Delta r$ in the $(J,a)$ plane stays positive everywhere---spin
    does not invert once the endpoints are pinned. The gradient is set by the
    endpoint angle ($J$), not by $a$.}
    \label{fig:cm-spin-fix}
\end{figure}


The verdict is independent of the endpoint symmetry: repeating both scans
with \emph{asymmetric} fixed endpoints ($r_A=10,r_B=6$) leaves
$\Delta r>0$ everywhere as well (spin $[+0.76,+2.64]$, conformal
$[+0.16,+3.24]$; Figs.~\ref{fig:asym-spin}--\ref{fig:asym-conf}),
consistent with the pointwise bound $n_t/n_\tau=E/f$.

\begin{table}[t]
    \centering
    \caption{Protocol dependence of plunge inversion. Same-launch
    comparison inverts under either rotation or a conformal factor;
    the genuine fixed-endpoint brachistochrone never inverts.}
    \label{tab:protocol}
    \begin{ruledtabular}
    \begin{tabular}{lcc}
        protocol & spin $a$ & conformal $A$ \\
        \colrule
        same launch & inverts (\S\ref{sec:inversion}) & inverts (\S\ref{sec:inversion}) \\
        fixed endpoints & no ($\Delta r>0$) & no ($n_t/n_\tau=E/f>1$) \\
    \end{tabular}
    \end{ruledtabular}
\end{table}

\subsection{The breathing indicatrix: a unified adiabatic theorem}
\label{sec:unified}
The Thakurta--Kerr and Vaidya corrections [Eqs.~\eqref{eq:psi-split},
\eqref{eq:vaidya-deltaphi}] are one statement, and the indicatrix ellipse of
Sec.~\ref{sec:indicatrix} is why. For a one-parameter family of rail indicatrices
with frozen integrand $F(r;\lambda)$ on the genus-two curve $y^2=S(r;\lambda)$
($\lambda=$ scale $a$, mass $m$, or conformal $A$), a slow drift
$\lambda=\lambda(\text{clock})$ gives
\begin{equation}
    \varphi=\varphi_0(r;\lambda)+\dot\lambda\!\int\!\partial_\lambda F\,
    (\text{clock})\,\dd r
    =\varphi_0+\dot\lambda\Big[\sum_{k<j} q_{kj}W_{kj}+(\text{elem.})\Big],
    \label{eq:unified}
\end{equation}
with $q_{kj}$ algebraic (via $\partial_\lambda F=N_\lambda/S^{3/2}$) and
$W_{kj}=\int(U_j\dd U_k-U_k\dd U_j)$ the genus-two polylogarithms of the frozen
curve. The transcendence is universal; the geometries differ only in which datum of
the ellipse breathes and which clock weights it,
$\partial_\lambda F=\partial_c F(\partial_\lambda c)+\partial_R F(\partial_\lambda
R)$: Thakurta--Kerr breathes the semi-axes $R$ (conformal scale, clocks $t,\tau$),
letter a third-kind dipole \emph{at infinity}; Vaidya breathes the centre $c_r$ (the
radial wind $\propto\sqrt{2m/r}$, clock the advanced time $v$), letter a $\dd\log$
\emph{at the horizon} $r=2m$. FLRW is the degenerate circle ($c=0$, isotropic $R$):
branches coincide, polylogs collapse to genus zero. The thread ``circle $\to$
radial-wind ellipse $\to$ angular-wind ellipse'' is the ladder of weight-two letters
the breathing switches on, on one hyperelliptic backbone.

\paragraph{Domain of validity.} The statement is a \emph{first-order} adiabatic
theorem resting on two \emph{exact} lemmas---the algebraic reduction
$\partial_\lambda F=N_\lambda/S^{3/2}$ with the unique rational $c_k$ from the
polynomial identity~\eqref{eq:reduction-identity} (nonsingular when $S$ has distinct
roots, i.e.\ off the separatrix), and the by-parts assembly~\eqref{eq:psi-split} via
the shuffle relation---so the closed form is rigorous, not fitted, where those hold.
Two honest caveats: it is $O(\dot\lambda)$ (the next order needs a further reduction);
and it \emph{fails at the turning points} $S(r)=0$, where
$\partial_\lambda F\sim N_\lambda/S^{3/2}$ diverges and $\dot\lambda\,\delta\varphi$ is
not uniformly valid---the closed forms hold on the smooth arcs and are joined across
turning points by the standard WKB connection formulae. On the separatrix the genus
drops and the reduction degenerates to the elliptic (Bloch--Wigner) case, a controlled
limit. Within these usual first-order WKB bounds the theorem is airtight.


% ======================================================================
\section{Conclusions}
\label{sec:concl}
% ----------------------------------------------------------------------
The non-stationary brachistochrone is a well-posed controlled-rail problem: the
invariant $-u\cdot W=\Ehat$ is legitimate via Pontryagin's principle, and the
selector $W$ is fixed canonically by the Killing--CKV--Kodama hierarchy, with
the Kodama energy conserved along the rail even when free fall does not conserve
it. FLRW is the degenerate base---freezing without branch splitting---while the
mass breaks the $t/\tau$ symmetry through its radial gravitational gradient (the
spin instead driving the rotational plunge inversion). On
the equatorial sector we obtained closed-form separatrices in Weierstrass
functions, continued through the ergosphere in Doran coordinates, and an analytic
cusp theorem for the $|J|<\Jc$ reflection. The plunge inversion resolves into a
three-way picture: rotation and the conformal factor invert the $t/\tau$
ordering under a same-launch comparison, whereas physical evaporation never does,
and \emph{no} inversion survives once both endpoints are fixed. Slow drift of the
metric yields closed-form first-order corrections whose transcendence is a
universal genus-two hyperelliptic dilogarithm on the frozen spectral curve; a
single breathing-indicatrix theorem unifies the three geometries, the weight-two
letter migrating from a third-kind dipole at infinity (Thakurta--Kerr) to a
$\dd\log$ at the horizon (Vaidya)---where it is the face of the
accretion/evaporation asymmetry and collapses on the separatrix to a tabulated
elliptic (Bloch--Wigner) dilogarithm. Outlook: the
exact Sultana--Dyer source
\cite{SultanaDyer2005}, the McVittie/Schwarzschild--de Sitter limits
\cite{McVittie1933,Kottler1918,Kaloper2010} with their photon surfaces
\cite{ClaudelViracVirbhadra2001}, escape maps, and freezing-surface dynamics.


\clearpage
% ======================================================================
\appendix
\section{Validation and consistency checks}
\label{app:validation}
We collect the numerical checks that underpin the analytic results. The
minimum principle is verified directly by perturbing the computed
brachistochrones and confirming that both $T_t$ and $T_\tau$ are minimized
(Fig.~\ref{fig:flrw-var} for the FLRW branch degeneracy, and
Fig.~\ref{fig:verifmin} for the two distinct Vaidya branches). Stationary-limit
residuals confirm that the non-stationary
Hamiltonian flow reproduces the closed-form Kerr expressions as $m'\to0$ and
$A\to1$ (Fig.~\ref{fig:vaidya-kerra0}); the Kodama energy is conserved to
$\sim10^{-16}$ along the rail (Fig.~\ref{fig:kodama}). The equatorial closed
forms are validated end-to-end by the three-method superposition
(ODE flow, quadrature, evaluated Weierstrass) through the ergosphere
(Fig.~\ref{fig:sep3}), and the fixed-endpoint no-inversion theorem is
checked over the asymmetric-endpoint grids (Figs.~\ref{fig:asym-spin}--%
\ref{fig:asym-conf}).

\subsection{Explicit closed forms}
\label{app:explicit}

All expressions below are at \emph{constant (frozen) conformal factor} $A$:
stationary Kerr with $E\to E_{\rm eff}=\Ehat/A$, $a\to s$, $J\to J/A$. They give
the instantaneous orbit shape on a Thakurta--Kerr slice; the time dependence
enters only through $E_{\rm eff}(\eta)$ (Sec.~\ref{sec:quasiconst}).

\paragraph{Equatorial $\tau$-branch separatrix in Weierstrass functions.} On the
$\tau$-separatrix $J=\Jc$ the doranTau factorization $\Delta-\Jc^2 w=f(r^2+c^2)$,
$c=a/E$, makes the
turning quartic $Q_4=r\big((E^2-1)r+2M\big)(r^2+c^2)$, and the factor $r^2+c^2$
divides $Q_4$ so that only the horizon poles survive. Writing
\begin{gather}
    A=\frac{Ma^2}{2E^2},\qquad B=\frac{a^2(E^2-1)}{12E^2},\notag\\
    g_2=\frac{a^4(E^2-1)^2}{12E^4}-\frac{a^2M^2}{E^2},\notag\\
    g_3=\frac{a^4\big[36E^2M^2(1-E^2)-a^2(E^2-1)^3\big]}{216E^6},
    \label{eq:wp-invariants}
\end{gather}
the uniformizing variable is $z(r)=\wp^{-1}\!\big(A/r+B;g_2,g_3\big)
=\int_0^r\dd r'/\sqrt{Q_4}$, the Abel map of the turning quartic $Q_4$: it is the
elliptic argument in which the radial motion linearizes, so that $r$ is recovered
by $A/r+B=\wp(z)$ and the physical range $r\in[r_{\min},\infty)$ maps to a
segment of $z$. With $r_\pm=M\pm\sqrt{M^2-a^2}$ the horizons,
\begin{equation}
\begin{aligned}
    \alpha_\pm&=\frac{r_\pm\big((E^2-1)r_\pm+2M\big)}{r_\pm-r_\mp},\\
    \lambda_\pm&=\frac{\alpha_\pm}{\sqrt{Q_4(r_\pm)}},\qquad v_\pm=z(r_\pm),
\end{aligned}
\end{equation}
$\wp(v_\pm)=A/r_\pm+B$, and
$\Lambda_0=(E^2-1)-\alpha_+/r_+-\alpha_-/r_-+2\lambda_+\zeta(v_+)+2\lambda_-\zeta(v_-)$,
the angle is linear in the uniformizing variable $z\equiv z(r)$ plus two
Weierstrass-$\sigma$ ratios,
\begin{equation}
\begin{split}
    \varphi(r)=\frac{a}{E}\Big[&\Lambda_0\,z(r)
      +\lambda_+\ln\frac{\sigma(z-v_+)}{\sigma(z+v_+)}\\
      &+\lambda_-\ln\frac{\sigma(z-v_-)}{\sigma(z+v_-)}\Big]+\text{const}.
\end{split}
    \label{eq:sep-phi}
\end{equation}
Because $Q_4(2M)=4M^2(4E^2M^2+a^2)>0$, Eq.~\eqref{eq:sep-phi} is regular
\emph{through} the ergosphere $\re=2M$; the only singularities are the
third-kind logarithms at the horizons. The conformal case is
$E\to E_{\rm eff}=\Ehat/A$ in every parameter.

\paragraph{Doran continuation through the horizon.} The Boyer--Lindquist form
$\varphi(r)$ winds logarithmically at $\rp$ (the standard BL artifact); adding the
Doran shift $\dd\varphi_{\rm shift}=2Mar\,\dd r/(\Delta\sqrt{C_3})$ with the cubic
$C_3=2Mr(r^2+a^2)$ regularizes it. This is a \emph{second} elliptic curve, in
clean Weierstrass form $r(z)=(2/M)\wp(z;g_2,g_3)$ with $g_2=-M^2a^2$, $g_3=0$
(quasi-lemniscatic), giving
$\varphi_{\rm shift}=\tfrac{a}{r_+-r_-}\big(B_+-B_-\big)$,
$B_\pm=2\zeta(v_\pm)z+\ln[\sigma(z-v_\pm)/\sigma(z+v_\pm)]$; the log divergences of
$\varphi$ and $\varphi_{\rm shift}$ cancel at $\rp$, the analytic content of Doran
regularity. Thus $\varphi_D=\varphi_{\rm BL}+\varphi_{\rm shift}$ is closed on two
elliptic curves.

\paragraph{General $\tau$-branch $\varphi(r)$: genus-2 Kleinian form.} Off the
separatrix, for \emph{arbitrary} $J$, the $\tau$ integrand
$\dd\varphi/\dd r=\mathcal K_\tau(r)/\sqrt{\mathcal S(r)}$ is genus~2, with the
sextic
\begin{equation}
    \mathcal S(r)=r(r-2M)\big[(E^2-1)r+2M\big]
      \big[r\Delta-J^2\big((E^2-1)r+2M\big)\big]
    \label{eq:tau-sextic}
\end{equation}
(six generically simple branch points, hyperelliptic) and the rational numerator
$\mathcal K_\tau=Jr(r-2M)\big[(E^2-1)r+2M\big]/\Delta$. Partial fractions split it
into a first-kind (holomorphic) part and third-kind poles at the horizons
$r_\pm=M\pm\sqrt{M^2-a^2}$,
\begin{equation}
    \varphi(r)=c_1\!\int\!\frac{r\,\dd r}{\sqrt{\mathcal S}}
      +c_0\!\int\!\frac{\dd r}{\sqrt{\mathcal S}}
      +\sum_{\pm}\alpha_\pm\!\int\!\frac{\dd r}{(r-r_\pm)\sqrt{\mathcal S}},
    \label{eq:tau-genus2}
\end{equation}
with closed-form coefficients (in $M,a,E,J$)
\begin{equation}
\begin{aligned}
    &c_1=J(E^2-1),\qquad c_0=2MJ,\\
    &\alpha_\pm=\mp\frac{Ja^2\big[M(E^2+1)\pm(E^2-1)\sqrt{M^2-a^2}\big]}
      {2\sqrt{M^2-a^2}}.
\end{aligned}
    \label{eq:tau-coeffs}
\end{equation}
The two holomorphic integrals are the Abel-map coordinates (arguments of the
Kleinian hyperelliptic $\sigma$), the horizon terms the third-kind
$\ln\sigma$-quotients---the exact genus-2 analogue of the $t$-branch $J_+^t$
[Eq.~\eqref{eq:t-genus2}]. We built $Y^2=\mathcal S$ in SageMath at $J=-0.9\Jc$
and computed its Riemann period matrix $\tau$; both Riemann relations hold
($\tau=\tau^\top$ to $7\times10^{-12}$, $\operatorname{Im}\tau\succ0$), and
Eq.~\eqref{eq:tau-genus2} reproduces the direct integral to $10^{-15}$. At
$|J|=\Jc$ the factor $f^2$ (from $\Delta-\Jc^2w=f(r^2+c^2)$) makes $r_e$ a
\emph{double} root of $\mathcal S$; the $(r-2M)$ in $\mathcal K_\tau$ cancels one
factor, dropping the radical to the quartic $Q_4$---genus~1, the Weierstrass
separatrix~\eqref{eq:sep-phi}. This double root is also why the retrograde
marginal $J=-\Jc$ reflects with a \emph{corner} (finite $\dd\varphi/\dd r$)
rather than a smooth grazing: only the prograde $+\Jc$, aided by co-rotation,
threads through.

\paragraph{$t$-branch closed form: generic integrand.} The coordinate-time
brachistochrone has $\dd\varphi/\dd r=\mathcal K(r)/\sqrt{\mathcal R(r)}$ with the
\emph{sextic} radical $\mathcal R=r\,Q_2(r)\,[(E^2-1)r+2M]$, genus~2, the quartic
being (no cubic term)
\begin{equation}
    Q_2(r)=E^2r^4+\big[E^2(a^2-J^2)+J^2\big]r^2
      +2M(J-a)\big[E^2(J-a)-2J\big]r+4M^2(J-a)^2,
    \label{eq:t-Q2}
\end{equation}
and a \emph{closed-form} numerator in the problem quantities
\begin{equation}
    \mathcal K(r)=\frac{r\,[(E^2-1)r+2M]\,\big(J(r-2M)+2Ma\big)}{\Delta(r)},
    \label{eq:t-K}
\end{equation}
holding for \emph{any} $(M,a,E,J)$: the third-kind poles sit at the two horizons
$r_\pm$ ($\Delta=0$) and the polynomial part carries the first-/second-kind
pieces (there is \emph{no} pole at $r=0$, cancelled by the leading factor $r$).
By the conformal transfer $n^2_{\rm TK}=n^2_{\rm Kerr}(E_{\rm eff})$
(Sec.~\ref{sec:thakurta}), Eq.~\eqref{eq:t-K} is also the Thakurta--Kerr
$t$-integrand under $E\to E_{\rm eff}=\Ehat/A$, $J\to J/A$, so the conformal
factor is \emph{not} a separate variable of $\varphi(r)$.

\paragraph{$t$-branch separatrix in Weierstrass functions.} At the saddle-node
$J_c^-$ [Eq.~\eqref{eq:Jc-minus}] the radical acquires a
\emph{double root} at the marginal-capture radius $r_*$,
$\mathcal R=(r-r_*)^2 Q_4(r)$, and
$\dd\varphi/\dd r=\mathcal G(r)/\sqrt{Q_4}$ becomes a third-kind elliptic
differential on the quartic $Q_4$ (which retains the root $r=0$). With the same
uniformization $z(r)=\wp^{-1}(A/r+B;g_2,g_3)$ (invariants of $Q_4$) the closed
form is
\begin{equation}
    \varphi(r)=\Lambda_0\,z(r)+\sum_{k}\lambda_k\,
      \ln\frac{\sigma(z-v_k)}{\sigma(z+v_k)},
    \label{eq:t-sep-phi}
\end{equation}
now with third-kind poles only at $c_k\in\{\rp,\rm,r_*\}$---three real ones, with
\emph{no} pole at $r=0$ and no complex pair, unlike the $\tau$-separatrix. Here
$v_k=\wp^{-1}(A/c_k+B)$, $\lambda_k=-A\alpha_k/(c_k^2\wp'(v_k))$, and
$\Lambda_0=\Lambda_{\rm poly}-\sum_k\alpha_k/c_k+\sum_k 2\zeta(v_k)\lambda_k$,
with $\alpha_k$ the residues of $\mathcal G$. For $M=1$, $a=0.9$, $E=1.2$
($r_*=3.514$, $J_c^-=-8.05$) Eq.~\eqref{eq:t-sep-phi} reproduces the Hamiltonian
flow to $2\times10^{-14}$ (Fig.~\ref{fig:t-sep}). So both separatrices are
genus~1 by different mechanisms; only non-separatrix $t$-orbits need genus-2
Kleinian functions.

\begin{figure}[ht]
    \centering
    \includegraphics[width=\columnwidth]{Immagini/fig_tbranch_separatrix.png}
    \caption{$t$-branch separatrix (the retrograde saddle-node $J_c^-$): the
    orbit spirals onto the marginal circle $r_*$. Top: Hamilton ODE vs the
    genus-1 closed form (the sextic radical has a double root, dropping to the
    quartic $Q_4$). Bottom: they agree to $\sim10^{-13}$; inset shows the double
    root of $R_6=(r-r_*)^2Q_4$.}
    \label{fig:t-sep}
\end{figure}

\paragraph{Genus-2 Kleinian form and ergosphere penetration.} Away from the
separatrices the sextic $\mathcal R$ keeps six \emph{simple} branch points
(here $\{-4.545,\,-1.722\pm1.793i,\,0,\,1.444,\,2.000\}$), so the curve
$y^2=\mathcal R(r)$ is a genuine genus-2 hyperelliptic Riemann surface and
$\varphi(r)$ closes only in \emph{Kleinian} (hyperelliptic $\sigma$) functions.
With the numerator~\eqref{eq:t-K} the partial-fraction split
$\mathcal K=c_1 r+c_0+\alpha_+/(r-\rp)+\alpha_-/(r-\rm)$ gives
\begin{equation}
    \varphi(r)=\underbrace{c_1\!\!\int\!\frac{r\,\dd r}{\sqrt{\mathcal R}}
      +c_0\!\!\int\!\frac{\dd r}{\sqrt{\mathcal R}}}_{\text{first kind}}
      +\underbrace{\sum_{\pm}\alpha_\pm\!\!\int\!
      \frac{\dd r}{(r-r_\pm)\sqrt{\mathcal R}}}_{\text{third kind}},
    \label{eq:t-genus2}
\end{equation}
the holomorphic pair being the Abel-map coordinates (arguments of the Kleinian
$\sigma$) and the horizon terms the third-kind $\ln\sigma$-quotients---the exact
genus-2 analogue of Eq.~\eqref{eq:t-sep-phi}. For $M=1$, $a=0.9$, $E=1.2$ we
built $y^2=\mathcal R$ in SageMath and computed its Riemann period matrix
$\tau$; both Riemann relations hold ($\tau=\tau^\top$ to $4\times10^{-12}$,
$\operatorname{Im}\tau\succ0$ with eigenvalues $\{0.61,1.08\}$), certifying the
homology basis. The single Abelian integral $\varphi(r)$ does not invert to a
one-$\theta$ quotient $r(\varphi)$ (that needs the 2D Jacobi inversion,
$r=\wp_{22}$ along the trajectory line in the Jacobian, à la
Hackmann--L\"ammerzahl); but Eq.~\eqref{eq:t-genus2} reproduces the Hamiltonian
flow to $2\times10^{-13}$. Physically these genus-2 orbits are exactly the ones
that \emph{penetrate the ergosphere}: the prograde threshold $J_+^t$
[Eq.~\eqref{eq:Jplus-t}] has its turning point \emph{on} $r_e=2M$ (simple root at
$r=2$), while for $J$ just inside the capture window $[J_c^-,J_+^t]$ the turning
point moves \emph{inside} the ergoregion and the closed form stays accurate
across and below $r_e$ (Fig.~\ref{fig:jpt-g2}). This is the \emph{asymmetric}
$t$-window $J\in[J_c^-,J_+^t]=[-8.05,3.43]$; the $\tau$-branch instead penetrates
over the \emph{symmetric} window $|J|<a/E$ (turning exactly on $r_e$ at
$|J|=a/E$), the two-sided asymmetry of the $t$-window being the frame-dragging
wind, absent from the $J^2$-even $\tau$ radial turning. Only the retrograde
$t$-separatrix $J_c^-$ (the saddle-node) never reaches the ergosphere---it
spirals onto $r_*=3.514>r_e$.

\begin{figure}[ht]
    \centering
    \includegraphics[width=\columnwidth]{Immagini/fig_jpt_genus2.png}
    \caption{Genus-2 prograde $t$-orbits and ergosphere penetration. Top: the
    threshold $J_+^t=3.43$ (grey) turns \emph{on} $r_e=2M$; a momentum just
    inside the capture window ($J=3.0$) penetrates the ergoregion (shaded),
    Hamilton ODE (black) and the genus-2 closed form~\eqref{eq:t-genus2}
    (dotted) coinciding across $r_e$. Bottom: the deviation stays
    $\sim2\times10^{-13}$ including \emph{inside} the ergosphere ($r<r_e$,
    shaded).}
    \label{fig:jpt-g2}
\end{figure}

The two $t$-window boundaries are themselves of different kinds
(Table~\ref{tab:penetration}): the prograde $J_+^t$ is a \emph{simple} root on
$r_e$ (the orbit reaches the ergosphere and crosses), whereas the retrograde
$J_c^-$ is a \emph{double} root at $r_*=3.514>r_e$---a logarithmic
$\varphi$-barrier onto which the orbit spirals with infinite winding, never
reaching $r_e$ even though the region $(r_e,r_*)$ is itself allowed
($\mathcal R>0$). At $J=J_c^-$ the closed form~\eqref{eq:t-genus2} accordingly
diverges as $r\to r_*$. Just inside the
window the mechanism is a \emph{saddle-node}: the double root splits into a
complex-conjugate pair ($r_*=3.514\to3.393\pm0.601i$ at $J=0.95\,J_c^-$), the
$\varphi$-barrier vanishes, and the retrograde orbit is captured---plunging
through $r_e$ to the horizon, with the genus-2 closed form tracking it to
$\sim10^{-9}$ (Fig.~\ref{fig:jcm-cap}).

\begin{figure}[ht]
    \centering
    \includegraphics[width=\columnwidth]{Immagini/fig_jcm_capture.png}
    \caption{Retrograde capture just inside the window ($J=0.95\,J_c^-=-7.65$).
    Top: the orbit plunges from $r_0=6$ through the ergosphere (shaded) to the
    horizon $r_+$ (Hamilton ODE, black; genus-2 closed form, dotted), whereas
    the marginal $J_c^-$ (grey) only spirals onto $r_*$. Bottom: the saddle-node
    in the complex-$r$ plane---the double root of $\mathcal R$ at $r_*$ (square)
    splits into a complex-conjugate pair (circles) as $J$ moves off $J_c^-$,
    removing the barrier and opening the plunge.}
    \label{fig:jcm-cap}
\end{figure}

\begin{table*}[t]
    \centering
    \small
    \caption{Ergosphere behaviour of the two branches. The $\tau$-branch is a
    \emph{trichotomy} that penetrates only at the single prograde $J=+J_c$
    ($J_c=a/E$): the reality wall $\sqrt{wf}$ confines it to $r\ge r_e$, and off
    $+J_c$ it reflects---smoothly for $|J|>J_c$, with a semicubic cusp at $r_e$
    for $|J|<J_c$ (the retrograde marginal $-J_c$ cusps too, counter-rotation
    being forbidden inside $r_e$). The $t$-branch is a \emph{dichotomy} that
    penetrates the whole interval $(J_c^-,J_+^t)$, frame dragging auto-tuning the
    crossing; it never cusps. Values for $M=1$, $a=0.9$, $E=1.2$; orbit shapes in
    Figs.~\ref{fig:atlas-tau} and~\ref{fig:atlas-t}.}
    \label{tab:penetration}
    \begin{tabular}{lccc}
    \hline
    branch & penetrates & smooth bounce (outside $r_e$) & cusp bounce (at $r_e$) \\
    \hline
    $\tau$ & only $J=+J_c=+0.75$ & $|J|>J_c$
      & $|J|<J_c$ \ (and $J=-J_c$) \\
    $t$ & $(J_c^-,J_+^t)=(-8.05,3.43)$ & $J\notin(J_c^-,J_+^t)$ & --- (none) \\
    \hline
    \end{tabular}
\end{table*}

\begin{figure}[t]
    \centering
    \includegraphics[width=\columnwidth]{Immagini/fig_atlas_tau.png}
    \caption{$\tau$-brachistochrones by regime (trichotomy). Each panel is one
    momentum relative to $\pm J_c$ (adaptive Hamiltonian flow): smooth bounce
    outside $r_e$ ($|J|>J_c$), the single penetrating $J=+J_c$ (crosses, spirals
    onto $r_+$), cusp bounce at $r_e$ ($|J|<J_c$), the retrograde marginal
    $J=-J_c$ (also a cusp), and a smooth bounce for $J<-J_c$. Grey disc:
    ergoregion.}
    \label{fig:atlas-tau}
\end{figure}

\begin{figure}[t]
    \centering
    \includegraphics[width=\columnwidth]{Immagini/fig_atlas_t.png}
    \caption{$t$-brachistochrones by regime (dichotomy). Relative to
    $J_c^\pm=(-8.05,3.43)$: smooth bounce ($J>J_c^+$), tangent contact at $r_e$
    ($J=J_c^+$), penetration ($J_c^-<J<J_c^+$), retrograde capture just inside
    the edge ($J=0.99\,J_c^-$, plunging to $r_+$), and a smooth bounce
    ($J<J_c^-$). Unlike $\tau$, a whole interval penetrates and no orbit cusps.}
    \label{fig:atlas-t}
\end{figure}

\paragraph{Stationary-limit conserved quantities.} At $A=1$ the branch flows admit
the closed first integrals used for validation,
\begin{equation}
    \mathcal K_\tau=J,\qquad
    \mathcal K_t=\frac{fJ+2Ms/r}{E},\qquad
    \mathcal K_\eta=\frac{fJ}{E},
    \label{eq:kerr-conserved}
\end{equation}
the equatorial Kerr doranTau/doranT constants.

\paragraph{Rotational inversion function.} The same-launch $t/\tau$ ordering
inverts where $\rho(r_*)^2=H_{t0}/H_{\tau0}$ with the explicit ratio of
centrifugal potentials
\begin{equation}
    \rho=\frac{A^2}{\Ehat}\Big(f+\frac{2Ms}{rJ}\Big),
    \label{eq:rho-explicit}
\end{equation}
$\rho<1$ giving $t$ deeper, $\rho>1$ giving $\tau$ deeper.

\paragraph{Conformal trichotomy.} With constant $A$ the separatrix momentum is
$\Jc=sA^2/\Ehat$, and above it the smooth periapsis is the outer root of
\begin{equation}
    \Delta(r)-\Big(\tfrac{J}{A}\Big)^2\Big[\big(\tfrac{\Ehat}{A}\big)^2-f(r)\Big]=0,
    \label{eq:periapsis-root}
\end{equation}
while $|J|<\Jc$ reflects at a cusp $r_{\min}=\re$ (as does the retrograde
marginal $-\Jc$, forced counter-rotation forbidding it inside $\re$), and only
the prograde $+\Jc$ crosses to $\rp$.

\paragraph{Closed toroidal averages.} The secular source $\langle S_1\rangle$ of
Eq.~\eqref{eq:S1} follows from the orbit-plane averages ($\cos\theta=\sin i\,
\sin\psi$, $k=1-J^2/L^2$)
\begin{equation}
    \Big\langle\cos^2\theta\Big\rangle=\tfrac{k}{2},\quad
    \Big\langle\tfrac{1}{\sin^2\theta}\Big\rangle=\tfrac{L}{|J|},\quad
    \big\langle p_\theta^2\big\rangle=L(L-|J|),
    \label{eq:averages}
\end{equation}
with $\langle\cos2\theta\,p_\theta^2\rangle=|J|(L-|J|)$.

\subsection{Reproducing the reductions}
\label{app:reproduce}
Every coefficient in the adiabatic corrections is fixed by a \emph{single}
polynomial identity the reader can solve directly; nothing is fitted. The depth-one
reduction $\partial_\lambda\varphi_0=\mathcal{A}/\sqrt S+\sum_{k=0}^4 c_k U_k$
[Eq.~\eqref{eq:dEphi0}; likewise the Vaidya mass reduction of
Sec.~\ref{sec:vaidya-adiabatic}], differentiated in $r$ and cleared of the radical
via $\partial_\lambda F=N_\lambda/S^{3/2}$, is equivalent to
\begin{equation}
    2\,\mathcal{A}'S-\mathcal{A}\,S'+2S\sum_{k=0}^4 c_k r^k=2N_\lambda,
    \qquad \mathcal{A}=\sum_{i=0}^5 a_i r^i,
    \label{eq:reduction-identity}
\end{equation}
a degree-ten polynomial identity. Matching the coefficients of $r^0,\dots,r^{10}$
gives a nonsingular $11\times11$ linear system in $\{a_0,\dots,a_5,c_0,\dots,c_4\}$;
one linear solve returns the $c_k$ rationally in the metric parameters ($S$ the
branch radical: $S$ on the $\tau$-branch, $R_6=r\,Q_2\,DE$ on the $t$-branch). Here
$N_\lambda=EJ\,r^4(r-2M)^2 DE$ for the $\tau$-branch energy reduction ($\Delta$
cancels), $N_\lambda=R_6\partial_E\mathcal K-\tfrac12\mathcal K\partial_E R_6$ for
the $t$-branch [with $Q_2,\mathcal K$ of Eqs.~\eqref{eq:t-Q2},~\eqref{eq:t-K}, again
polynomial of degree seven], and $N_m=S\partial_m K-\tfrac12K\partial_m S$ for
Vaidya. The $t$-branch clock is entirely on-curve [Eq.~\eqref{eq:clock-t}]: its
second-kind vector $b$ feeds the $W_{kj}$ of $R_6$, its horizon part
$R_\Delta/\Delta$ the third-kind $\ln\sigma$-quotients; no double cover. Three
reference solutions (polynomial residual $0$ exactly):
\begin{table*}[t]
    \centering
    \caption{Reference reduction coefficients $c_k$ from
    Eq.~\eqref{eq:reduction-identity}, reproducible by one linear solve; the script
    \texttt{reproduce\_reductions.py} performs all three.}
    \label{tab:reduction-ck}
    \begin{ruledtabular}
    \begin{tabular}{lllllll}
        reduction & parameters & $c_0$ & $c_1$ & $c_2$ & $c_3$ & $c_4$\\
        \hline
        TK $\tau$ & $M{=}1,a{=}0.9,E{=}1.2,J{=}2.5$
            & $-1.1575$ & $1.86406$ & $-0.57499$ & $0.02429$ & $0.07117$\\
        TK $t$ & $M{=}1,a{=}0.9,E{=}1.2,J{=}2.5$
            & $-0.39084$ & $1.72654$ & $-0.19058$ & $-0.14503$ & $-0.01676$\\
        Vaidya & $M{=}1,E{=}1.4,J{=}2.5$
            & $-0.87153$ & $-0.35970$ & $1.03538$ & $0.21292$ & $-0.18939$\\
    \end{tabular}
    \end{ruledtabular}
\end{table*}

Given the $c_k$, the weight-two assembly~\eqref{eq:psi-split} is mechanical: the
clock supplies $b$ ($\tau$-branch $\eta=U_3-2M U_2$, $b=(0,0,-2M,1,0)$), the weights
are $Q_{kj}=c_kb_j-c_jb_k$, the letters $W_{kj}=\int(U_k\dd U_j-U_j\dd U_k)$, and
$\mathcal{I}_j=\int\mathcal{A}r^j/S\,\dd r$ integrate to logarithms; Vaidya follows
identically with the advanced-time clock [Eq.~\eqref{eq:vaidya-full}], the horizon
dilogarithm closing by one integration by parts [Eq.~\eqref{eq:horizon-dilog}].
Table~\ref{tab:script-map} maps each result to the verifying script in the
repository~\cite{Rosignoli2026Code}; the first entry runs the whole chain in one
file---from $(M,a,E,J)$ through the $11\times11$ solve and the building blocks to the
assembled $\delta\varphi$, matching the direct integral to $4\times10^{-15}$.
\begin{table*}[t]
    \centering
    \caption{Result-to-script map with achieved numerical residual.}
    \label{tab:script-map}
    \begin{ruledtabular}
    \begin{tabular}{lll}
        result & script & residual\\
        \hline
        \emph{end-to-end} $(M,a,E,J)\!\to\!\delta\varphi$ & \texttt{pipeline\_completa\_deltaphi.py} & $4{\times}10^{-15}$\\
        TK $\tau$ reduction $c_k$, $\psi$ & \texttt{kerr\_psi\_explicit\_verified.py} & $5{\times}10^{-14}$\\
        TK $t$-branch clock \& assembly & \texttt{kerr\_tbranch\_psi\_assembly.py} & $9{\times}10^{-13}$\\
        quasi-periods $\eta$; $L$ (Fay) & \texttt{kerr\_quasiperiods\_bel.sage} & $10^{-12}$; $5{\times}10^{-8}$\\
        Vaidya reduction $c_k^m$ & \texttt{vaidya\_dMF\_reduction.py} & $10^{-15}$\\
        Vaidya clock; $\delta\varphi$~\eqref{eq:vaidya-full} & \texttt{vaidya\_fully\_explicit.py} & $2{\times}10^{-13}$\\
        Vaidya $\tau$-branch assembly & \texttt{vaidya\_tau\_assembly.py} & $10^{-14}$\\
        horizon dilog $\mathcal{D}_k$ & \texttt{vaidya\_horizon\_dilog.py} & $10^{-12}$\\
        separatrix $U_0(\sigma,\zeta)$; $\tau$ & \texttt{vaidya\_ell\_dilog\_match.py} & $10^{-14}$\\
        accr./evap.\ split & \texttt{vaidya\_asymmetry.py} & $10^{-15}$\\
    \end{tabular}
    \end{ruledtabular}
\end{table*}

\subsection{Consistency residuals}
\label{app:residuals}
The closed forms and reductions above were checked symbolically and numerically.
The Weierstrass separatrix~\eqref{eq:sep-phi} reproduces the direct quadrature to
$|\Delta\varphi|<2.8\times10^{-10}$ through the ergosphere, with the $\wp$-identity
satisfied to $9\times10^{-28}$ and the antiderivative to $30$ digits; the Doran
shift reduces to $1.1\times10^{-16}$. The non-autonomous Randers
data~\eqref{eq:randers-data} match the Kerr null Fermat metric to
$\sim10^{-14}$. The equatorial Thakurta--Kerr flow reproduces the $A=1$ Kerr
constants~\eqref{eq:kerr-conserved} to ten digits at $r=4,6,9$ ($a=0.9$, $E=1.2$),
and the $s\to0$ limit the Thakurta--Schwarzschild flow to $10^{-12}$; the $3+1$
Hamiltonian~\eqref{eq:H3d} reduces to the equatorial one exactly, and conserves
$L^2$ to $6\times10^{-13}$ at $a=0$. The Kodama energy is held to
$|{-}u\cdot K-\Ehat|<7\times10^{-16}$ along the rail, and the static
inversion/trichotomy colormap agrees with the analytic
contours~\eqref{eq:periapsis-root} to a relative $7\times10^{-12}$.

\begin{figure}[tbp]
    \centering
    \includegraphics[width=\columnwidth]{Immagini/fig_colormap_conforme_estremi_fissi.png}
    \caption{Conformal scan at fixed symmetric endpoints (via
    $E_{\rm eff}=\hat E/A$): $\Delta r$ in the $(J,A)$ plane is positive
    everywhere, matching the closed-form theorem $n_t/n_\tau=E/f>1$.}
    \label{fig:cm-conf-fix}
\end{figure}


\begin{figure}[tbp]
    \centering
    \includegraphics[width=\columnwidth]{Immagini/fig_colormap_spin_asimmetrico.png}
    \caption{Robustness of the fixed-endpoint no-inversion result under
    \emph{asymmetric} endpoints ($r_A=10,r_B=6$), spin scan (Kerr flow):
    $\Delta r>0$ throughout the $(J,a)$ plane.}
    \label{fig:asym-spin}
\end{figure}

\begin{figure}[tbp]
    \centering
    \includegraphics[width=\columnwidth]{Immagini/fig_colormap_conforme_asimmetrico.png}
    \caption{Same robustness check, conformal scan (via $E_{\rm eff}=\hat E/A$):
    $\Delta r>0$ throughout the $(J,A)$ plane for asymmetric endpoints. The
    blank wedge is the region where the $\tau$-turning point rises above
    $r_B$ (no dipping orbit reaches $B$).}
    \label{fig:asym-conf}
\end{figure}

\clearpage

\begin{figure}[tbp]
    \centering
    \includegraphics[width=\columnwidth]{Immagini/fig_vaidya_kerr_a0.png}
    \caption{Stationary-limit check ($m'=0$): numerical orbit vs closed-form
    Kerr $a=0$.}
    \label{fig:vaidya-kerra0}
\end{figure}


\begin{figure}[tbp]
    \centering
    \includegraphics[width=\columnwidth]{Immagini/fig_verifica_minimo_brachi.png}
    \caption{Each branch minimizes its own travel time; $t$ and $\tau$ are
    distinct curves.}
    \label{fig:verifmin}
\end{figure}

\begin{figure}[tbp]
    \centering
    \includegraphics[width=\columnwidth]{Immagini/fig_vaidya_bounce.png}
    \caption{Bounce through periapsis in the $v$-parametrization: the
    Euler--Lagrange flow is regular at $r'=0$ and integrates through the turning
    point, reaching $r_{\min}$ exactly where the $r$-parametrization diverges.}
    \label{fig:vaidya-bounce}
\end{figure}

\begin{figure}[tbp]
    \centering
    \includegraphics[width=\columnwidth]{Immagini/fig_vaidya_timing.png}
    \caption{Timing window: $r_{\min}$ depends on the arrival epoch $v_1$. The
    orbit--horizon race has opposite winners in absolute ($r_{\min}$) and relative
    ($r_{\min}/2m$) terms for accretion vs evaporation.}
    \label{fig:vaidya-timing}
\end{figure}

\clearpage

% ======================================================================
\begin{acknowledgments}
The author thanks the University of Pavia for providing a stimulating academic
environment. The author received no external funding for this work. The AI
assistants GPT~5.5 (OpenAI) and Claude Opus~4.8 (Anthropic) were used to assist with
symbolic derivations, numerical implementation, and manuscript preparation; all
results were verified independently by the author, who takes full responsibility
for the content of this work. The closed-form and hyperelliptic computations relied
on open-source software, and the author gratefully acknowledges their developers:
\textsc{SageMath}~\cite{SageMath} (The Sage Developers) and the
\textsc{abelfunctions} library~\cite{abelfunctions} (C.~Swierczewski and
collaborators) for Riemann surfaces, Abel maps, and Riemann theta functions with
derivatives; and \textsc{SymPy}~\cite{SymPy}, \textsc{NumPy}~\cite{NumPy},
\textsc{SciPy}~\cite{SciPy}, and \textsc{Matplotlib}~\cite{Matplotlib} for the
symbolic reductions, numerical integration, and figures.
\end{acknowledgments}

\section*{Data availability}
The complete numerical and symbolic codebase underlying this work is openly
available under the MIT license in the GitHub repository
``NonStationary\-Brachi\-sto\-chrone''~\cite{Rosignoli2026Code}:
\begin{center}
\url{https://github.com/ImanetorX98/NonStationaryBrachistochrone}
\end{center}
The repository includes the trajectory integrators, the Pontryagin Hamiltonian
constructions, the Poisson-bracket closure and quasi-invariant scans
(Sec.~\ref{sec:quasiconst}), the Weierstrass separatrix and Doran-continuation
scripts, the adiabatic-reduction and hyperelliptic-dilogarithm pipelines mapped in
Table~\ref{tab:script-map}, the fixed-endpoint BVP and colormap pipelines, and the
figure-generation code. In particular the self-contained script
\texttt{reproduce\_reductions.py} regenerates the coefficients of
Table~\ref{tab:reduction-ck} from the polynomial
identity~\eqref{eq:reduction-identity} alone (one linear solve, residual zero). It
also retains exploratory scripts and negative results,
documenting which ansatz classes failed the held-out trajectory and phase-space
residual tests reported in Sec.~\ref{sec:quasiconst} and
App.~\ref{app:validation}.

\clearpage
\bibliographystyle{apsrev4-2}
\bibliography{refs}

\end{document}
